% !TEX program = lualatex
\documentclass[10pt, leqno]{ltjsarticle}
% \documentclass[10pt]{ltjsarticle} % 日本語をそのまま扱えるクラス
% \usepackage[margin=28mm]{geometry}
\usepackage{amsmath,amssymb,mathtools}
\usepackage{amsthm}
\usepackage{tikz}
\usetikzlibrary{arrows.meta,positioning,calc}

\usepackage[
  top=15mm,     % 上余白
  bottom=15mm,  % 下余白
  left=10mm,    % 左余白
  right=10mm,    % 右余白
  landscape,
  margin=10mm
]{geometry}
% \usepackage[landscape,margin=20mm]{geometry} % 余白はお好みで
\usepackage{paracol}
\columnratio{0.5,0.5}             % 左右を 1:1 に
\setlength{\columnseprule}{0.4pt} % 真ん中の縦線

\usepackage{titlesec}
\titleformat{\paragraph}
  {\normalfont\bfseries}
  {}      % ← 番号を空に
  {0pt}
  {}

% ---- 記号・マクロ ----
\newcommand{\xto}{\mathrel{\xrightarrow{}}}
\newcommand{\into}{\hookrightarrow}
\newcommand{\epi}{\twoheadrightarrow}
\newcommand{\pto}{\rightrightarrows}
\newcommand{\cls}[1]{\bigl(#1\bigr)}
\newcommand{\redjoin}{\wedge}
\newcommand{\shrink}{\phi}
\newcommand{\Ex}{E}        % suspension functor E
\newcommand{\Ox}{\Omega}   % loop functor Omega
\newcommand{\id}{\mathrm{id}}
% --- Toda bracket helpers ---
\newcommand{\toda}[1]{\{\,#1\,\}}              % {a,b,c}
\newcommand{\todai}[2][]{\{\,#2\,\}_{#1}}      % {a,b,c}_i   (#1 は添字, 空なら無し) 
\newcommand{\gen}[1]{\{\,#1\,\}}       % ⟨generator⟩ 表記を { } で統一
\newcommand{\Ker}{\operatorname{Ker}}
\newcommand{\im}{\operatorname{Im}}
\newcommand{\hs}[2]{\bigl[#1,#2\bigr]}   % 一般のホモトピー集合
\newcommand{\pis}[2]{\pi_{#1}(#2)}       % 左が球面のときのホモトピー群
\newcommand{\pit}[2]{\pi_{#1}^{#2}}
\newcommand{\ext}[1]{\overline{#1}}     % extension
\newcommand{\coext}[1]{\widetilde{#1}}  % coextension
\newcommand{\z}{\mathbb{Z}}

% ---- theorem-like environments ----
\numberwithin{equation}{section}
\theoremstyle{plain}
\newtheorem{theorem}{Theorem}[section]
\newtheorem{lemma}[theorem]{Lemma}
\newtheorem{proposition}[theorem]{Proposition}
\newtheorem{corollary}[theorem]{Corollary}
\theoremstyle{definition}
\newtheorem{definition}[theorem]{Definition}
\newtheorem{example}[theorem]{Example}
\newtheorem{notation}[theorem]{Notation}
\newtheorem{assumption}[theorem]{Assumption}
\theoremstyle{remark}
\newtheorem{remark}[]{Remark}

% \title{PART 1 \quad CHAPTER I\\[2mm] General Composition and Secondary Composition}
\author{}
\date{}

\begin{document}

\setcounter{section}{5}

\section*{CHAPTER V.\quad Auxiliary Calculation of \(\pi_{n+k}(S^n;2)\) for \(1\le k\le 7\)}

\begin{paracol}{2}


本章では, (4.4) の完全列を使うことによって, \\
\(1\le k\le 7\) に対する \(\pi_{n+k}(S^n)\) の 2-成分に関して,\\
\([7][8][14][16][20]\) で知られている結果について述べる. \\
以下, \(\mathbb Z\) を整数全体から成る群, 
\(\mathbb Z_r\) を法 \(r\) の整数から成る群とする. \\
さらに, \(\gen{\alpha}\subset \pi_i(S^n)\) を, 
\(\alpha\in \pi_i(S^n)\) により生成される巡回部分群とする. \\
(4.1) から, 次が成り立つ：
\begin{equation}\label{5.1}
\pi^{1}_{i}=0\ (i>1),\qquad
\pi^{n}_{i}=0\ (i<n),\qquad
G_k=(G_k;2)=0\ (k<0).
\end{equation}
\[
\pi^{n}_{n}=\gen{\iota_n}\cong \mathbb Z,\qquad
G_0=\gen{\iota}\cong \mathbb Z. \quad 
\iota=E^{\infty}\iota_n.
\]

\paragraph*{i) The Groups \(\pi^{n}_{n+1}\).}
(5.1) より, \(\pi^{1}_{2}=0\) である. Prop 4.4 から,
\(
H:\pi^{2}_{3}\longrightarrow \pi^{3}_{3}
\)
は同型である. \\
\(H(\eta_2)=\iota_3\) を満たす類 \(\eta_2\in \pi^{2}_{3}\) をとると, 
\(\eta_2\) は \(\pi^{2}_{3}\cong\mathbb Z\) を生成する. \\
\(n\ge2\) に対し, \(\eta_n=E^{\,n-2}\eta_2\), 
さらに, \(\eta=E^{\infty}\eta_2\) と記す. \\
Prop 2.7 により,
\(H(\Delta(\iota_5))=\pm H[\iota_2,\iota_2]=\pm 2\iota_3\) である. \\
\(H:\pi^{2}_{3}\to\pi^{3}_{3}\) が同型なので, 
\[
\Delta(\iota_5)=\pm 2\,\eta_2 .
\]
完全列 (4.4)
\[
\pi^{5}_{5}\xrightarrow{\ \Delta\ }\pi^{2}_{3}\xrightarrow{\ E\ }\pi^{3}_{4}\xrightarrow{\ \ }\pi_{4}^{5}=0
\]
を考える. ここで, \(E\) は全射であり, 
その kernel は \(2\eta_2\) により生成される. \\
したがって,
\(\pi^{3}_{4}=\gen{\eta_3}\cong\mathbb Z_2\) であり, (4.5) から,
\[
\pi^{n}_{n+1}=\gen{\eta_n}\cong\mathbb Z_2, \ (n\ge3), \
(G_1;2)=\gen{\eta}\cong\mathbb Z_2.
\]
以上より, 次を得る. 

\switchcolumn

\noindent
\textbf{\cite{Freudenthal1937}
H.~Freudenthal,}
\newblock Über die Klassen der Sphärenabbildungen,
\newblock Compos. Math., {\bf 5} (1937).

\noindent
\textbf{\cite{Hopf1935}
H.~Hopf,}
\newblock Über die Abbildungen von Sphären auf Sphären niedriger Dimension,
\newblock Fund. Math., {\bf 25} (1935).

\noindent
\textbf{\cite{Serre1953a}
J.-P.~Serre,}
\newblock Cohomologie modulo 2 des complexes d'Eilenberg-MacLane,
\newblock Comm. Math. Helv., {\bf 27} (1953), 198--231.

\noindent
\textbf{\cite{Serre1953c}
J.-P.~Serre,}
\newblock Quelques calculs de groupes d'homotopie,
\newblock Comptes Rendus Paris, {\bf 236} (1953), 2475--2477.

\noindent
\textbf{\cite{Toda1952}
H.~Toda,}
\newblock Generalized Whitehead products and homotopy groups of spheres,
\newblock Jour. of Inst. Poly. Osaka City Univ., {\bf 3} (1952), 43--82.

\noindent
\textbf{(4.1).}\
\(\pis{i}{S^1} = 0\ (i>1),\
\pis{i}{S^n} = 0\ (i < n),\
\pis{n}{S^n} = \z\gen{\iota_n}. \)

\noindent
\textbf{Prop4.4.}\
\(H(\alpha) = \pm 2\iota_{2m-1}\ \Longrightarrow\
\pit{i-1}{m-1}\oplus\pit{i}{2m-1}\xrightarrow{\cong}\pit{i}{m},\
(\beta,\gamma)\mapsto E\beta+\alpha\circ\gamma.\)\\
\(0=\pit{2}{1}\to\pit{3}{2}
\xrightarrow{H}{\cong}\pit{3}{3}
\xrightarrow{\Delta}\pit{1}{1}
\xrightarrow{E}\pit{2}{2}.\)

\noindent
\textbf{Prop2.7.}\
\(H\Delta(\iota_{2n+1}) = \pm2\iota_{2m-1},\ (n:\text{even})\).

\noindent
\textbf{(4.5).}\
\(E^{m-n}:\pit{n+k}{n}\xrightarrow{\cong}\pit{m+k}{m},\
(n\ge k+2). \)


\end{paracol}

\begin{paracol}{2}

\begin{proposition}\label{prop:5.1}
% Prop5.1.\
\(\pi^{2}_{3}=\gen{\eta_2}\cong\mathbb Z,\quad
\pi^{n}_{n+1}=\gen{\eta_n}\cong\mathbb Z_2\ (n\ge3),\quad
(G_1;2)=\gen{\eta}\cong\mathbb Z_2.\) \\
ここで, 
\(H(\eta_2)=\iota_3,\ \Delta(\iota_5)=\pm 2\,\eta_2\) が成り立つ. 
\end{proposition}

(5.1) より,
\(\pi^{1}_{i-1}=0\) (\(i\ge3\)) なので, Prop 4.4 から, 次が成り立つ. 
\begin{equation}\label{eq:5.2}
% (5.2).\
\text{合成 } \alpha\longmapsto \eta_2\circ \alpha \text{ は, 同型 } 
\pi^{3}_{i}\xrightarrow{\cong}\pi^{2}_{i}\ (i\ge3) 
\text{ を定める. }
\end{equation}

以下のLemを, 後で用いる. 

\begin{lemma}\label{lem:5.2}
% \noindent
% \textbf{Lem5.2.}\
\(\alpha\in\pi_{i}(S^{3})\) に対し,
\(2\alpha=0\) と仮定する. \\
このとき,
\(\todai[1]{\eta_3,\,2\iota_4,\,E\alpha}\) 
の任意の元 \(\beta\) に対して,
\[
H(\beta)=E^{2}\alpha,\qquad
2\beta=\eta_3\circ E\alpha\circ \eta_{i+1}
\]
が成り立つ. そのような \(\beta\) は,
\(\pi^{3}_{i+2}\) に属し, \(\Delta(E^{2}\alpha)=0\) である. 
\end{lemma}

\switchcolumn

\noindent
\textbf{Prop4.4.}\
\(\alpha\in \pit{2n-1}{n},\ H(\alpha)=\pm\iota_{2n-1}\) のとき,\\
\(\pit{i-1}{n-1}\oplus\pit{i}{2n-1}\to\pit{i}{n},\ 
(\beta, \gamma)\mapsto E\beta+ \alpha\circ\gamma\)
は, \({}^\forall i\) で同型.


\end{paracol}

\begin{paracol}{2}

\noindent
\textbf{Proof.}\ 
Lem 4.5 より, \(2\iota_3\circ \alpha=2\alpha=0\). \\
Toda bracket 
\(\todai[1]{\eta_3,\,2E\iota_3,\,E\alpha}\) に, 
Prop 2.6 を適用すると, 
\[
H(\beta)\in \Delta^{-1}(\eta_2\circ 2\iota_3)\circ E^{2}\alpha
= \Delta^{-1}(2\eta_2)\circ E^{2}\alpha
= \pm \iota_5\circ E^{2}\alpha
= E^{2}\alpha.
\]
また,
\[
\begin{align*}
2\beta
& \in  2\todai[1]{\eta_3,\,2\iota_4,\,E\alpha}\\
& =  
\todai[1]{\eta_3,\,2\iota_4,\,E\alpha}\circ 2\iota_{i+2}\\
& =  
\eta_3\circ E\toda{2\iota_3,\alpha,2\iota_{i}} 
&& \text{by Prop1.4}\\
& \subset 
\eta_3\circ -\,\toda{2\iota_{4},E\alpha,2\iota_{i+1}}
&& \text{by Prop1.3}\\
& \ni  
\eta_3\circ -(E\alpha\circ\eta_{i+1})
&& \text{by Cor3.7}\\
& =  \eta_3\circ E\alpha\circ \eta_{i+1}.
\end{align*}
\]
さらに, 
\(\eta_3\circ -\todai[1]{2\,\iota_{4},E\alpha,2\,\iota_{i+1}}\) は,
\[
\eta_3\circ \pi_{i+2}(S^{4})\circ 2\iota_{i+2} + \eta_3\circ2\iota_4\circ E\pi_{i+2}(S^4) = \eta_3\circ\pi_{i+2}(S^4)\circ2\iota_{i+2}
\] 
の coset である. \\
\(\gamma\in\pi_{i+2}(S^4)\) に対して, (2.1) により, 
\(E(\eta_3\circ\gamma\circ2\iota_{i+2})=\eta_4\circ2E\gamma
=2\eta_4\circ E\gamma=0 .\)\\
Lem 4.5 から, \(\eta_3\circ\gamma\circ2\iota_{i+2}=0 .\)\\
よって, \(\eta_3\circ\pi_{i+2}(S^4)\circ2\iota_{i+2}=0\) なので, \\
\(\eta_3\circ - \todai[1]{2\iota_4,E\alpha,2\iota_{i+1}}\)
は, ただ1つの元から成る. \\
以上より, \(2\beta=\eta_3\circ E\alpha\circ\eta_{i+1}\) である. \\
明らかに, 
\(4\beta=0,\ \beta\in\pi^{3}_{i+2}\), であり,
\(\Delta(E^{2}\alpha)=\Delta(H(\beta))=0\) である. 
\qed

\switchcolumn

\noindent
\textbf{Lem4.5.}\
\(n=2,4,8\) のとき, 
\(E:\pis{i-1}{S^{n-1}}\to\pis{i}{S^n}\) は単射.\\ 
\(r\iota_{n-1}\circ\beta = r\beta,\ \beta\in\pis{i-1}{S^{n-1}}\).

\noindent
\textbf{Prop2.6.}\
\(H\todai[1]{E\alpha, E\beta, E\gamma} 
= -\Delta^{-1}(\alpha\circ\beta)\circ E^2\gamma.\)

\noindent
\textbf{Prop1.4.}\
\(\todai[n]{\alpha, E^n\beta, E^n\gamma}\circ E^{n+1}\delta
= (-1)^{n+1}(\alpha\circ E^n\toda{\beta, \gamma, \delta})\).

\noindent
\textbf{Prop1.3.}\
\(-E\todai[n]{\alpha, E^n\beta, E^n\gamma}
\subset \todai[n+1]{E\alpha, E^{n+1}\beta, E^{n+1}\gamma}\).

\noindent
\textbf{Cor3.7.}\
\(\todai[1]{r\iota_{m+1}, E\beta, r\iota_{m+k+1}}
\ni E\beta\circ\eta_{m+k+1}\)\ 
if \ \(r\equiv 2\ \pmod 4\).

\noindent
\(S^4\xleftarrow{2\iota_4}
S^4\xleftarrow{E\alpha}
S^{i+1}\xleftarrow{2\iota_{i+1}}
S^{i+1}\).

\noindent
\(\eta_3\circ 2\iota_4 = 0.\)

\noindent
\textbf{(2.1).}\
\(k(\alpha\circ\beta) = \alpha\circ k\beta,\
k(\alpha\circ E\beta) = (k\alpha)\circ E\beta.\)

\noindent
\textbf{Lem4.5.}\
\(E: \pis{i-1}{S^3}\to\pis{i}{S^4}\) は単射.



\end{paracol}

\begin{paracol}{2}

\(\eta_3\in\pi_4(S^3)\) は, 上のLemの条件 \(2\iota_3\circ\eta_3=0\) を満たす. \\
\(\nu'\in\todai[1]{\eta_3,2\iota_4,\eta_4}\) を考えると, 
次を得る. 
\begin{equation}\label{eq:5.3}
% (5.3).\
\nu'\in\pi_6^3,\ H(\nu')=\eta_5,\ 2\nu'=\eta_3\circ\eta_4\circ\eta_5 .
\end{equation}


\paragraph*{ii) The Group \(\pi^{n}_{n+2}\).}
\(\eta_n^{2}:=\eta_n\circ \eta_{n+1}\) および,
\(\eta^{2}:=\eta\circ\eta\) と記す. このとき,
\[
E\,\eta_n^{2}=\eta_{n+1}^{2},\qquad E^{\infty}\eta_n^{2}=\eta^{2}.
\]

\switchcolumn

\end{paracol}

\begin{paracol}{2}

\begin{proposition}\label{prop:5.3}
% Prop5.3.\
\(\pi^{n}_{n+2}=\gen{\eta_n^{2}}\cong \mathbb Z_2\ (n\ge2).\) \
\((G_2;2)=\gen{\eta^{2}}\cong \mathbb Z_2 .\)
\end{proposition}

\noindent
\textbf{Proof.}\ 
 (5.2) より, \(\pi_{4}^{2}=\gen{\eta_2^{2}}\cong \mathbb Z_2 .\)
完全列 (4.4)
\[
\pi_{6}^{3}\xrightarrow{\ H}\pi_{6}^{5}\xrightarrow{\ \Delta\ }\pi_{4}^{2}
\xrightarrow{\ E\ }\pi_{5}^{3}\xrightarrow{\ H\ }\pi^{5}_{5}
\xrightarrow{\ \Delta\ }\pi_{3}^{2}
\]
を考える. Prop 5.1 から,
\(\Delta:\pi^{5}_{5}\to\pi_{3}^{2}\) は単射である. \\
したがって, \(H(\pi_{5}^{3})=0\) で, \(E\) は全射. \\
さらに, (5.3) とProp 5.1 より,
\(H:\pi^{3}_{6}\to\pi^{5}_{6}\) は全射である. \\
ゆえに, \(\Delta(\pi^{5}_{6})=0\) であり, 
\(E\) は単射となる. \\
よって, \(E:\pi_4^2\xrightarrow{}\pi_5^3\) は同型となり, \(n=3\) の場合が示せた. \\
次に, 
\(\pi_5^3\xrightarrow{E}\pi_6^4\xrightarrow{}\pi_6^7\)
において, (5.1) により,
\(\pit{6}{7}=0\) なので,
\(E\) は全射である. \\
Prop4.4 により, \(E\) は単射なので, 
\(E:\pit{5}{3}\to\pit{6}{4}\) は同型であるので, \(n=4\) の場合が示せた. \\
\(n>4\) と, \((G_2;2)\) については, (4.5) から分かる. 
\qed

\switchcolumn

\noindent
\textbf{(5.2).}\
\(\text{合成 } \alpha\longmapsto \eta_2\circ \alpha \text{ は, 同型 } 
\pi^{3}_{i}\xrightarrow{\cong}\pi^{2}_{i}\ (i\ge3) 
\text{ を定める. }\)

\noindent
\textbf{Prop5.1.}\
\(\pi^{2}_{3}=\gen{\eta_2}\cong\mathbb Z,\quad
\pi^{n}_{n+1}=\gen{\eta_n}\cong\mathbb Z_2\ (n\ge3),\quad
(G_1;2)=\gen{\eta}\cong\mathbb Z_2.\) \\
ここで, 
\(H(\eta_2)=\iota_3,\ \Delta(\iota_5)=\pm 2\,\eta_3\) が成り立つ. 

\noindent
\textbf{(5.3).}\
\(\nu'\in\pi_6^3,\ H(\nu')=\eta_5,\ 2\nu'=\eta_3\circ\eta_4\circ\eta_5 .\)

\noindent
\textbf{Prop4.4.}\
\(H(\alpha) = \pm 2\iota_{2m-1}\ \Longrightarrow\
\pit{i-1}{m-1}\oplus\pit{i}{2m-1}\xrightarrow{\cong}\pit{i}{m},\
(\beta,\gamma)\mapsto E\beta+\alpha\circ\gamma.\)\\
\(0=\pit{2}{1}\to\pit{3}{2}
\xrightarrow{H}{\cong}\pit{3}{3}
\xrightarrow{\Delta}\pit{1}{1}
\xrightarrow{E}\pit{2}{2}.\)

\noindent
\textbf{(4.5).}\
\(E^{m-n}:\pit{n+k}{n}\xrightarrow{\cong}\pit{m+k}{m},\
(n\ge k+2). \)


\end{paracol}

\begin{paracol}{2}

\bigskip
\paragraph*{iii). The Group \(\pi^{n}_{\,n+3}\).}

\begin{lemma}[]
\label{lem:5.4}
% Lem5.4.\
ある元 \(\nu_{4}\in \pi^{4}_{7}\) が存在して, 次を満たす. 
\[
H(\nu_{4})=\iota_7,\qquad 2E \nu_{4}=E^{2}\nu' .
\]
\end{lemma}

\noindent
\textbf{Proof.}\ 
Th 3.6 を, \(\alpha=\eta_{2}\) に適用すると, 
\(\eta_{2}\) は条件 \((1-(-1)^{1})E^{4}\eta_{2}=2\eta_{6}=0\) を満たす. 
よって, ある元 \(\alpha^{\ast}\in\pis{7}{S^4}\) が存在し, \\
\(\beta\circ E^{t}\eta_{2}=0\) を満たす
任意の \(\beta\in\pi_{\,t+2}(S^{\,m})\) に対して,
\[
E^{2}\beta\circ E^t\alpha^{\ast}\in (-1)^m \toda{E^{m}\eta_{2},\,E^{3}\beta,\,E^{t+3}\eta_{2}}_{3}
\]
となる. 
とくに, \(\beta=2\iota_{3}\) とし, \(t=1\) の場合,  (2.1) より
\(2\iota_3\circ E\eta_{2}=2E\eta_{2}=2\eta_3=0\). \\
したがって,
\[
E^{2}(2\iota_{3})\circ E\alpha^{\ast}=2E\alpha^{\ast}
\]
は, \(-\todai[3]{\eta_{5},\,2\iota_{6},\,\eta_{6}}\) に含まれる. 
ここで,
\(n\ge3,\ t\le n-2\) のとき, 
\begin{equation}\label{eq:5.4}
% (5.4).\
\todai[t]{\eta_{n},\,2\iota_{n+1},\,\eta_{n+1}} = \{E^{\,n-3}\nu',-E^{\,n-3}\nu'\}.
\end{equation}

\noindent
まず, \(t\ge1\) の場合を考える.  
(4.7) , Prop 5.3 と (5.3) より, 
上の Toda bracket は,
\[
\eta_n\circ \pi^{n+1}_{n+3} + \pi^{n}_{n+2}\circ \eta_{n+2} 
= \gen{\eta_n\circ\eta_{n+1}\circ\eta_{n+2}}
= \gen{2E^{n-3}\nu'}
\]
の coset である. 

\switchcolumn

\noindent
\textbf{Th3.6.}\
$\alpha\in\pis{n+k}{S^{n}}$, $h\ge0$ が \eqref{3.1} を満たすとする.\\ 
さらに, $k\le 2n-2$ と仮定する.
このとき,
\(\alpha^{*}\in\pis{2n+2k+1}{S^{2n}}\)
が存在して, \\
\(\beta\circ \Ex^{\ell}\alpha=0\)
を満たす任意の $\beta\in\pis{n+\ell}{S^{n}}$ に対して, \\
\(\Ex^{n}\beta \circ \Ex^{t}\alpha^{*}
\in (-1)^{km+kt+t}\,\todai[n+k]{\Ex^{m}\alpha,\ 
  \Ex^{n+k}\beta,\ \Ex^{n+k+t}\alpha}\)\\
\(+ (-1)^{kn+h+t+1}\,\todai[h+t]{\Ex^{n}\beta,\ 
  \Ex^{n+t}\alpha,\ (1-(-1)^k)\Ex^{n+k+t}\alpha}\ .\)\\
さらに,
\((1-(-1)^k)\Ex^{n+h+1}\alpha=0\), 特に \(k\) が偶数のとき, \\
\(
\Ex^{n}\beta\circ \Ex^{t}\alpha^{*}\ \in\
(-1)^{km+kt+t}\,\todai[n+k]{\Ex^{m}\alpha,\ \Ex^{n+k}\beta,\ \Ex^{n+k+t}\alpha}\ .
\)

\noindent
\textbf{(2.1).}\
\(k(\alpha\circ\beta) = \alpha\circ k\beta,\
k(\alpha\circ E\beta) = (k\alpha)\circ E\beta.\)

\noindent
\textbf{(4.7).}\
\(n\ge1, \ \alpha\in\pis{k+n}{S^m; 2},\ \gamma\in\pis{i}{S^j; 2}
    \ \text{のとき, }\)\\
\(\alpha\circ E^n\pis{j}{S^k} = \alpha\circ E^n \pit{j}{k}
  \subset \pit{j+n}{m},
  \ \pis{j}{S^k}\circ\gamma = \pit{j}{k}\circ\gamma \subset\pit{i}{k}\ . \)\\
\(\text{さらに, }\ \alpha\circ E^n\beta=0,\ \beta\circ\gamma = 0 \ \text{のとき, }\)\\
\(\todai[n]{\alpha, E^n\beta, E^n\gamma}\ \text{は,}\
  \alpha\circ E^n\pit{i+1}{k} + \pit{j+n+1}{m}\circ E^{n+1}\gamma\
  \text{の coset で, }\pit{i+n+1}{m} \ \text{の部分集合である. }\)

\noindent
\textbf{(5.3).}\
\(\nu'\in\pi_6^3,\ H(\nu')=\eta_5,\ 2\nu'=\eta_3\circ\eta_4\circ\eta_5 .\)




\end{paracol}

\begin{paracol}{2}

\noindent
\(\nu'\) の定義, Prop 1.3, (1.15) から, 
\begin{align*}
E^{\,n-3}\nu'\in E^{\,n-3}\todai[1]{\eta_3,\,2\iota_{4},\,\eta_4}
& \subset(-1)^{n-3}\todai[n-2]{\eta_n, 2\iota_{n+1}, \eta_{n+1}}\\
& \subset(-1)^{n-3}\todai[t]{\eta_n, 2\iota_{n+1}, \eta_{n+1}}
\end{align*}
が従うので, \(t\ge1\) で (5.4) が成立する. \\
(3.2) より, \(\pi_{\,n+3}(S^{n+1})=E\,\pi_{\,n+2}(S^{n})\) である. \\
したがって,
\(\todai[1]{\eta_n,2\iota_{n+1},\eta_{n+1}}\) と
\(\toda{\eta_n,2\iota_{n+1},\eta_{n+1}}\)
は, 同じ次の群の coset である. \\
\(
\eta_n\circ E\pi_{n+2}(S^n) + \pi_{n+2}(S^n)\circ\eta_{n+2} .
\)\\
(1.15) から, これらの Toda bracket は一致する. \\
これから, \(t=0\) のとき, (5.4) が成り立つことが分かる. 
\qed

今, \(2E\alpha^{\ast}=\pm E^{2}\nu'\) である. \\
\(H:\pi^{3}_{6}\to \pi^{5}_{6}\) は,
\(\nu'\) を \(\pi_6^5\) の生成元 \(\eta_{5}\) に写すので, \(\nu'\) は 2 で割れない. \\
一方, \(2E\alpha^{\ast}=\pm E^{2}\nu'\) より,
\(E^{2}\nu'\) は 2 で割れる. \\
(4.8) より,
\(H(\alpha^{\ast})=(2s+1)\iota_{7}\), (\(s\in\mathbb Z\)). \\
そこで,
\(H[\iota_4,\iota_4]=(-1)^u 2\iota_7\) とし, 
\(\nu_4\) を次で定義する. 
(\([\iota_4,\iota_4]\) は Whitehead product)
\[
\nu_{4}:=
\begin{cases}
\alpha^{\ast}-(-1)^{u}\,s[\iota_4,\iota_4] 
& (\text{if } 2E\alpha^{\ast}=E^{2}\nu'),\\[1mm]
-\alpha^{\ast}+(-1)^{u}(s+1)[\iota_4,\iota_4] 
& (\text{if } 2E\alpha^{\ast}=-E^{2}\nu').
\end{cases}
\]
このとき, \(\nu_4\) がLem5.4 の条件を満たすことは簡単に確かめられる. 

\begin{lemma}\label{lem:5.5}
% Lem5.5.\
$Lem 5.4$ の元 \(\nu_{4}\) に対し, 
Toda bracket 
\(\todai[3]{\eta_{m+2},\,E^{3}\beta,\,\eta_{t+5}}\) は,\\
\(
E^{2}\beta\circ E^{t}\nu_{4}\text{または,} -E^{2}\beta\circ E^{t}\nu_{4}
\)
を含む. \\
ここで, \(\beta\in\pi_{t+2}(S^{m})\) は,
\(\beta\circ \eta_{t+2}=0\), \(t>0\) を満たす元である. 
\end{lemma}

\noindent
\textbf{Proof.}\ 
前Lemの証明より, \\
\(E^{2}\beta\circ E^{t}\alpha^{\ast}\) あるいは,
\(-E^{2}\beta\circ E^{t}\alpha^{\ast}\) 
が, 当該 Toda bracket に含まれる. \\
\(E[\iota_4,\iota_4]=0\) なので, 定義より,
\(E^{t}\nu_{4}=\pm E^{t}\alpha^{\ast}\) である. 
\qed

\[
\nu_n:=E^{\,n-4}\nu_{4},\ (n\ge 4),\qquad \nu:=E^{\infty}\nu_{4}
\]
\[
\eta^{3}_{n}:=\eta_{n}\circ\eta_{n+1}\circ\eta_{n+2},\ (n\ge2),\qquad
\eta^{3}:=\eta\circ\eta\circ\eta .
\]
と記す.
Lem 5.4 とProp 5.3 より,
\begin{equation}\label{eq:5.5}
% (5.5).\
n\ge5\ \text{に対して, } 2\nu_n=E^{\,n-3}\nu',\qquad
4\nu_{n}=\eta_n^{3},\quad 4\nu=\eta^{3}.
\end{equation}
\noindent
Prop 4.4 とLem 5.4 から,
\begin{equation}\label{eq:5.6}
% (5.6).\
(\alpha,\beta)\ \longmapsto\ E\alpha+\nu_{4}\circ\beta
\quad\text{は,}\quad
\pi^{\,3}_{i-1}\oplus \pi^{\,7}_{i}\ \xrightarrow{\ \cong\ }\ \pi_{\,i}^{4}
\ \text{の同型を与える. }
\end{equation}

\switchcolumn

\noindent
\textbf{(Def).}\
\(\nu'\in\todai[1]{\eta_3,2\iota_4,\eta_4}\).

\noindent
\textbf{Prop1.3.}\
\(-E\todai[n]{\alpha, E^n\beta, E^n\gamma}
\subset \todai[n+1]{E\alpha, E^{n+1}\beta, E^{n+1}\gamma}\).

\noindent
\textbf{(1.15).}\
\(\{\alpha,\ \Ex^n\beta,\ \Ex^n\gamma\}_n
\ \subset\
\bigl\{\alpha,\ \Ex^{m}(E^{n-m}\beta),\ E^{m}(E^{n-m}\gamma)\bigr\}_m .\)

\noindent
\textbf{(3.2).}\
\(\Ex: \pis{i}{S^m} \longrightarrow \pis{i+1}{S^{m+1}}
\quad\text{は, } i<2m-1 \text{で同型, } i=2m-1 \text{で全射である. }\)

\noindent
\textbf{(5.4).}\
\(\todai[t]{\eta_{n},\,2\iota_{n+1},\,\eta_{n+1}} = \{E^{\,n-3}\nu',-E^{\,n-3}\nu'\}.\)

\noindent
\textbf{(4.8).}\
\( n \ \text{を偶数とする. } \)\\
\(\Delta(\iota_{2n-1})=[\iota_{n-1},\iota_{n-1}]\ne0 \ \text{のとき, }\)\\
\(\pit{2n-1}{n}
=E\,\pit{2n-2}{\,n-1}\ \oplus\ \Delta\,\pit{2n+1}{\,2n+1}
\cong E\,\pit{2n-2}{\,n-1}\ \oplus\ \mathbb{Z},\ \pit{2n}{\,n+1} = E^{2}\,\pit{2n-2}{\,n-1} 
\cong E\,\pit{2n-2}{\,n}.\)\\
\(\Delta(\iota_{2n-1})=0 \ \text{のとき, }\)\\
\(E^{2}:\pit{2n-2}{\,n-1}\to\pit{2n}{\,n+1}\ \text{は単射で, }\
\pit{2n}{\,n+1}/E^{2}\pit{2n-2}{\,n-1}\cong\mathbb{Z}_2 . \)\\
\(\text{さらに,}\ \alpha\in\pit{2n-1}{n} \ \text{に対して, 次が成り立つ. }\)\\
\(E\alpha\in E^{2}\pit{2n-2}{\,n-1} \ \Leftrightarrow \
H(\alpha)\in 2\,\pit{2n-1}{\,2n-1} .\)

\noindent
\textbf{Lem5.4.}\
ある元 \(\nu_{4}\in \pi^{4}_{7}\) が存在して, 次を満たす. \\
\(H(\nu_{4})=\iota_7,\ 2E \nu_{4}=E^{2}\nu' .\)

\noindent
\textbf{Prop5.3.}\
\(\pi^{n}_{n+2}=\gen{\eta_n^{2}}\cong \mathbb Z_2\ (n\ge2).\) \
\((G_2;2)=\gen{\eta^{2}}\cong \mathbb Z_2 .\)

\noindent
\textbf{Prop4.4.}\
\(\alpha\in \pit{2n-1}{n},\ H(\alpha)=\pm\iota_{2n-1}\) のとき,\\
\(\pit{i-1}{n-1}\oplus\pit{i}{2n-1}\to\pit{i}{n},\ 
(\beta, \gamma)\mapsto E\beta+ \alpha\circ\gamma\)
は, \({}^\forall i\) で同型.





\end{paracol}

\begin{paracol}{2}

\begin{proposition}\label{prop:5.6}
% Prop5.6.\
\ \\
\(\pi^{2}_{5}=\gen{\eta_{2}^{3}}\cong \mathbb Z_{2},\) \\
\(\pi^{3}_{6}=\gen{\nu'}\cong \mathbb Z_{4},\) \\
\(\pi^{4}_{7}=\gen{\nu_{4}}\oplus \gen{E\nu'}\cong \mathbb Z\oplus\mathbb Z_{4},\) \\
\(\pi^{n}_{n+3}=\gen{\nu_{n}}\cong \mathbb Z_{8}\ (n\ge5),\)\qquad
\((G_{3};2)=\gen{\nu}\cong \mathbb Z_{8}.\)
\end{proposition}

\noindent
\textbf{Proof.}\ 
(5.2), Prop 5.3 より,
\(\pi^{2}_{5}=\gen{\eta_{2}^{3}}\cong\mathbb Z_{2}\). 
完全列 (4.4)
\[
\pi^{3}_{7}\xrightarrow{H}\pi^{5}_{7}\xrightarrow{\Delta}\pi^{2}_{5}
\xrightarrow{E}\pi^{3}_{6}\xrightarrow{H}\pi^{5}_{6}
\]
を考える. Prop 2.2 とProp 5.3 から,
\(
H(\nu'\circ \eta_{6})=H(\nu')\circ\eta_{6}
=\eta_{5}\circ \eta_{6}=\eta^{2}_{5}.
\)\\
また, Prop 5.3 より,
\(H:\pi_{7}^{3}\to \pi_{7}^{5}\) は全射で, 
\(E\) は単射. \\
ゆえに, \(\eta^{3}_{3}=E\eta^{3}_{2}\) の位数は 2. 
(5.3) から, \(\nu'\) の位数は 4である. \\
\(\pi_6^5=\mathbb Z_2\gen{\eta_5}\) なので, 
\(\nu'\) は \(\pi_6^3\) の生成元である. 
したがって, \(\pi^{3}_{6}=\gen{\nu'}\cong\mathbb Z_{4}\). \\
(5.6) より, \(\pi^{\,4}_{7}=\gen{\nu_{4}}\oplus \gen{E\nu'}
\cong \mathbb Z\oplus\mathbb Z_{4}\). \\
さらに, (4.8) と (5.6)? から, 
\(\pi^{5}_{8}/E^{2}\pi^{3}_{6}\cong \mathbb Z_{2}\) で, 
\(E^{2}:\pi^{3}_{6}\to \pi^{5}_{8}\) は単射. \\
(5.5) より,
\(\nu_5\) の位数は 8 であり, \(\pi_8^5\) の生成元である. \\
(4.5) により, 
\(\pi_{n+3}^n=\mathbb Z_8\gen{\nu_n},\ \) (\(n\ge5\)) , 
\((G_3;2)=\mathbb Z_8\gen{\nu}\) である. 
\qed

証明の中で, 次を得た. 
\begin{equation}\label{eq:5.7}
% (5.7).\
H(\nu'\circ\eta_6)=\eta_5^2 .
\end{equation}

\(E:\pi_7^4\rightarrow\pi_8^5\) の kernel が,
\(2\nu_4-E\nu'\) により生成されることは, 簡単に確かめられる. \\
これと, (4.4) の完全性から, 次を得る. 
\begin{equation}\label{eq:5.8}
% (5.8).\
\Delta(\iota_9)=\pm(2\nu_4-E\nu')  \ (=\pm[\iota_4, \iota_4]) .
\end{equation}

\paragraph*{iv). The Group \(\pi^{n}_{n+4}\).}

(5.2), Prop5.6 により, 
\(\pi_6^2=\mathbb Z_4\gen{\eta_2\circ\nu'}\) である. 

\begin{lemma}\label{lemma:5.7}
% Lem5.7.\
\(\alpha\in\pi_i(S^3)\) に対して,
\(E^2\alpha\in2\iota_5\circ\pi_{i+2}(S^5)\) ならば, 
\(E(\eta_2\circ\alpha)=0\) である. \\
特に, \(E(\eta_2\circ\nu')=0, \ \Delta(\nu_5)=\pm(\eta_2\circ\nu')\) である. 
\end{lemma}

\noindent
\textbf{Proof.}\ 
\(E^{2}(\eta_2\circ \alpha)
=\eta_4\circ E^2\alpha\in\eta_4\circ2\iota_5\circ\pi_{i+2}(S^5)=0.\)\\
従って, Lem4.5 から,
\(E(\eta_2\circ \alpha)=0\) である.\\
Lem5.4 により, \(E^2\nu'=2\nu_5=2\iota_5\circ E\nu_4\) なので, 
\(E(\eta_2\circ \nu')=0\) である. \\
列 (4.4) の完全性により,
\(\pi^{2}_{6}=\Delta(\pi^{5}_{8})=\gen{\Delta(\nu_5)} .\)\\
よって, \(\Delta(\nu_5)=\pm(\eta_2\circ\nu')\) である. 
\qed

\switchcolumn

\noindent
\textbf{(5.2).}\
\(\text{合成 } \alpha\longmapsto \eta_2\circ \alpha \text{ は, 同型 } 
\pi^{3}_{i}\xrightarrow{\cong}\pi^{2}_{i}\ (i\ge3) 
\text{ を定める. }\)

\noindent
\textbf{Prop5.3.}\
\(\pi^{n}_{n+2}=\gen{\eta_n^{2}}\cong \mathbb Z_2\ (n\ge2).\) \
\((G_2;2)=\gen{\eta^{2}}\cong \mathbb Z_2 .\)

\noindent
\textbf{(5.3).}\
\(\nu'\in\pi_6^3,\ H(\nu')=\eta_5,\ 2\nu'=\eta_3\circ\eta_4\circ\eta_5 .\)

\noindent
\textbf{(5.6).}\
\((\alpha,\beta)\ \longmapsto\ E\alpha+\nu_{4}\circ\beta
\quad\text{は,}\quad
\pi^{\,3}_{i-1}\oplus \pi^{\,7}_{i}\ \xrightarrow{\ \cong\ }\ \pi_{\,i}^{4}
\ \text{の同型を与える. }\)

\noindent
\textbf{(4.8).}\
\( n \ \text{を偶数とする. } \)\\
\(\Delta(\iota_{2n-1})=[\iota_{n-1},\iota_{n-1}]\ne0 \ \text{のとき, }\)\\
\(\pit{2n-1}{n}
=E\,\pit{2n-2}{\,n-1}\ \oplus\ \Delta\,\pit{2n+1}{\,2n+1}
\cong E\,\pit{2n-2}{\,n-1}\ \oplus\ \mathbb{Z},\ \pit{2n}{\,n+1} = E^{2}\,\pit{2n-2}{\,n-1} 
\cong E\,\pit{2n-2}{\,n}.\)\\
\(\Delta(\iota_{2n-1})=0 \ \text{のとき, }\)\\
\(E^{2}:\pit{2n-2}{\,n-1}\to\pit{2n}{\,n+1}\ \text{は単射で, }\
\pit{2n}{\,n+1}/E^{2}\pit{2n-2}{\,n-1}\cong\mathbb{Z}_2 . \)\\
\(\text{さらに,}\ \alpha\in\pit{2n-1}{n} \ \text{に対して, 次が成り立つ. }\)\\
\(E\alpha\in E^{2}\pit{2n-2}{\,n-1} \ \Leftrightarrow \
H(\alpha)\in 2\,\pit{2n-1}{\,2n-1} .\)

\noindent
\textbf{(5.5).}\
\(n\ge5\ \text{に対して, } 2\nu_n=E^{\,n-3}\nu',\
4\nu_{n}=\eta_n^{3},\quad 4\nu=\eta^{3}.\)

\noindent
\textbf{(4.5).}\
\(E^{m-n}:\pit{n+k}{n}\xrightarrow{\cong}\pit{m+k}{m},\
(n\ge k+2). \)

\noindent
\textbf{Lem4.5.}\
\(n=2,4,8\) のとき, 
\(E:\pis{i-1}{S^n-1}\to\pis{i}{S^n}\) は単射.\\ 
\(r\iota_{n-1}\circ\beta = r\beta,\ \beta\in\pis{i-1}{S^{n-1}}\).

\noindent
\textbf{Lem5.4.}\
ある元 \(\nu_{4}\in \pi^{4}_{7}\) が存在して, 次を満たす. \\
\(H(\nu_{4})=\iota_7,\ 2E \nu_{4}=E^{2}\nu' .\)


\end{paracol}

\begin{paracol}{2}


\(E\pi_6^2=\gen{E(\eta_2\circ\nu')}=0\) なので, 
(4.4) から, 完全列 
\(0\to\pi_7^3\xrightarrow H \pi_7^5\) を得る. \\
(5.7) とProp5.3 により, \(H\) は同型で, 
\(\pi_7^3=\mathbb Z_2\gen{\nu'\circ\eta_6}\) である. \\
これで, 次のPropの最初の2つが証明できた. 

\begin{proposition}\label{prop:5.8}
\ \\
% Prop5.8.\
\(\pi^{2}_{6}=\mathbb Z_4\gen{\eta_{2}\circ\nu'},\) \\
\(\pi^{3}_{7}=\mathbb Z_2\gen{\nu'\circ\eta_6},\) \\
\(\pi^{4}_{8}=\mathbb Z_2\gen{\nu_{4}\circ\eta_7}\oplus \mathbb Z_2\gen{E\nu'\circ\eta_7},\) \\
\(\pi^{5}_{9}=\mathbb Z_2\gen{\nu_{5}\circ\eta_8},\) \\
\(\pi^{n}_{n+4}=0\ (n\ge6),\) \
\((G_{4};2)=0.\)
\end{proposition}

\noindent
\textbf{Proof.}\ 
\(\pi_8^4\) に対する結果は, (5.6) とProp5.1 から従う. 
(4.4) の完全列
\[
\pi^{9}_{10}\xrightarrow{\ \Delta\ }\pi^{9}_{8}\xrightarrow{\ E\ }\pi^{5}_{9}
\xrightarrow{\ H }\pi^{9}_{9}\xrightarrow{\ \Delta\ }\pi^{4}_{7}
\]
を考える. (5.8) とProp5.6 により, \(\Delta:\pi_9^9\to\pi_7^4\) は単射である. \\
よって, 完全性から, \(E\) は全射となることが分かる. \\
\(E\) の kernel は次で生成される. 
\[
\begin{align*}
\Delta(\eta_9)
& =  \Delta(\iota_9)\circ \eta_7
&& \text{by Prop2.5}\\
& =  (2\nu_4-E\nu')\circ \eta_7
&& \text{by (5.8)}\\
& =  \nu_4\circ (2\eta_7)-E\nu'\circ \eta_7
 =  E\nu'\circ \eta_7.
\end{align*}
\]
これから, \(\pi_9^5=E\pi_8^4=\mathbb Z_2\gen{\nu_5\circ\eta_8}\) である. 
次に, (4.4) の完全列
\[
\pi^{11}_{11}\xrightarrow{\ \Delta\ }\pi^{5}_{9}\xrightarrow{\ E } \pi^{6}_{10}
\xrightarrow{\ H } \pi^{11}_{10}=0
\]
を考える. 
\(E\) は全射で, \(\pi_{10}^6\) は \(\nu_6\circ\eta_9\) で生成される. 
ここで, 次の関係式を証明する. 
\begin{equation}\label{eq:5.9}
% (5.9).\
\eta_3\circ \nu_4= \nu'\circ \eta_6,\qquad
\eta_n\circ \nu_{n+1}=0\ (n\ge5),\qquad
\nu_n\circ \eta_{n+3}=0\ (n\ge6)
\end{equation}

まず, \(\eta_3\circ\nu_4\) を考える. 
\(
2E(\eta_3\circ\nu_4)
=2(\eta_4\circ\nu_5)
=2\eta_4\circ\nu_5=0
\)
なので, \\
Lem4.5 から, \(2(\eta_3\circ\nu_4)=0 .\)
よって, \(\eta_3\circ\nu_4\in\pi_7^3 .\)\\
\(\pi_7^3\) は \(\nu'\circ\eta_6\) で生成されるので, 
\(\eta_3\circ\nu_4=x\nu'\circ\eta_6 \ (x = 0 \text{ or } 1) .\)
\[
\begin{align*}
H(\eta_3\circ \nu_4)
& =  E(\eta_2\wedge\eta_2)\circ H(\nu_4)
&& \text{by Prop2.2}\\
& =  \pm (\eta_5\circ \eta_6)\circ \iota_7
&& \text{by Prop3.1, Lem5.4}\\
& =  \eta_5^{2}
&& \text{by Prop5.3}\\
& =  H\big(x\nu'\circ \eta_6\big)
=x\eta_5^2 
&& \text{by (5.7)}.
\end{align*}
\]
これより, \(x=1\), 
\(\eta_3\circ\nu_4=\nu'\circ\eta_6\) である. 

\switchcolumn

\noindent
\textbf{(5.7).}\
\(H(\nu'\circ\eta_6)=\eta_5^2 .\)

\noindent
\textbf{Prop5.3.}\
\(\pi^{n}_{n+2}=\gen{\eta_n^{2}}\cong \mathbb Z_2\ (n\ge2).\) \
\((G_2;2)=\gen{\eta^{2}}\cong \mathbb Z_2 .\)

\noindent
\textbf{(5.6).}\
\((\alpha,\beta)\ \longmapsto\ E\alpha+\nu_{4}\circ\beta
\quad\text{は,}\quad
\pi^{\,3}_{i-1}\oplus \pi^{\,7}_{i}\ \xrightarrow{\ \cong\ }\ \pi_{\,i}^{4}
\ \text{の同型を与える. }\)

\noindent
\textbf{Prop5.1.}\
\(\pi^{2}_{3}=\gen{\eta_2}\cong\mathbb Z,\quad
\pi^{n}_{n+1}=\gen{\eta_n}\cong\mathbb Z_2\ (n\ge3),\quad
(G_1;2)=\gen{\eta}\cong\mathbb Z_2.\) \\
ここで, 
\(H(\eta_2)=\iota_3,\ \Delta(\iota_5)=\pm 2\,\eta_3\) が成り立つ. 

\noindent
\textbf{(5.8).}\
\(\Delta(\iota_9)=\pm(2\nu_4-E\nu')  \ (=\pm[\iota_4, \iota_4]) .\)

\noindent
\textbf{Prop5.6.}\
\(\pi^{2}_{5}=\gen{\eta_{2}^{3}}\cong \mathbb Z_{2},\) \\
\(\pi^{3}_{6}=\gen{\nu'}\cong \mathbb Z_{4},\) \\
\(\pi^{4}_{7}=\gen{\nu_{4}}\oplus \gen{E\nu'}\cong \mathbb Z\oplus\mathbb Z_{4},\) \\
\(\pi^{n}_{n+3}=\gen{\nu_{n}}\cong \mathbb Z_{8}\ (n\ge5),\)\qquad
\((G_{3};2)=\gen{\nu}\cong \mathbb Z_{8}.\)

\noindent
\textbf{Prop2.5.}\
$\alpha\in\pi_{i+1}(S^{2m+1}), \ \beta\in \pi_{j-1}(S^{i-1})$ 対して, 
\(
\Delta\alpha\ \circ\ \beta\ \subset\ \Delta(\alpha\circ E^2\beta).
\)\\
特に, 任意の $\gamma$ に対し,
\(
\Delta(E^2\gamma)\ \ni\ \pm\,[\,\iota_m,\ \iota_m\,]\circ\gamma
\)
が成り立つ. \\
ここで $[\, ,\,]$ は Whitehead product を表す. \\
上の「$\subset$」および「$\ni$」は, $m$ が奇数, 
あるいは $j<3m-1$, \\
または, 関係する元が \(2\)-primary 成分にあるときには, 
等号「$=$」に置き換えられる. 

\noindent
\textbf{Lem4.5.}\
\(n=2,4,8\) のとき, 
\(E:\pis{i-1}{S^n-1}\to\pis{i}{S^n}\) は単射.\\ 
\(r\iota_{n-1}\circ\beta = r\beta,\ \beta\in\pis{i-1}{S^{n-1}}\).

\noindent
\textbf{Prop2.2.}\
$K,L$ を有限 cell complex とする. \\
$\alpha\in[EK, S^{m+1}], \ \beta\in[L,K], \ \gamma\in\pi_r(S^r)$ 
に対して, 次が成り立つ：\\
\(
H(\alpha\circ E\beta)\ =\ H(\alpha)\circ E\beta,
\)\\
\(
H(E\gamma\circ\alpha)\ =\ E(\gamma\wedge\gamma)\circ H(\alpha).
\)

\noindent
\textbf{Prop3.1.}\
\(
\alpha \wedge \beta
= (-1)^{(p+k)h} E^{q}\alpha \circ E^{p+k}\beta
= (-1)^{ph} E^{p}\beta \circ E^{q+h}\alpha .
\)

\noindent
\textbf{Lem5.4.}\
ある元 \(\nu_{4}\in \pi^{4}_{7}\) が存在して, 次を満たす. \\
\(H(\nu_{4})=\iota_7,\ 2E \nu_{4}=E^{2}\nu' .\)

\noindent
\textbf{Prop5.3.}\
\(\pi^{n}_{n+2}=\gen{\eta_n^{2}}\cong \mathbb Z_2\ (n\ge2).\) \
\((G_2;2)=\gen{\eta^{2}}\cong \mathbb Z_2 .\)

\noindent
\textbf{(5.7).}\
\(H(\nu'\circ\eta_6)=\eta_5^2 .\)




\end{paracol}

\begin{paracol}{2}

次に, (5.5) により,
\(\eta_5\circ\nu_6=E^2\nu'\circ\eta_8=2\nu_5\circ\eta_8=0 .\)\\
よって, \(n\geq5\) に対して, \(\eta_n\circ\nu_{n+1}=0\) である.\\ 
Prop3.1 により, \(\eta_2\wedge\nu_4=\eta_6\circ\nu_7=\nu_6\circ\eta_9 .\)\\
よって, \(\nu_6\circ\eta_9=\eta_6\circ\nu_7=0, \)
\(\nu_n\circ\eta_{n+3}=0,\ (n\geq6)\) である. \\
これで, (5.9) が全て証明できた. 

最後の結果により, \(\pi_{10}^6=\gen{\nu_6\circ\eta_9}=0 .\)\\
(4.5) により, \(\pi_{n+4}^n=(G_4;2)=0, \ (n\geq6)\) である. \\
以上で, Prop5.8 の証明が完了した. 
\qed

証明の中で, \(E:\pi_9^5\to\pi_{10}^6\) が自明だったので, 
\(\Delta:\pi_{11}^{11}\to\pi_9^5\) が全射となった. \\
\(\pi_9^5\) が \(\nu_5\circ\eta_8\) で生成されるので, 次を得る. 
\begin{equation}\label{eq:5.10}
% (5.10).\
\Delta(\iota_{11})=\nu_5\circ\eta_8 .
\end{equation}

また, 証明の中で次の関係も得られた. 
\begin{equation}\label{eq:5.11}
% (5.11).\
\Delta(\eta_9)=E\nu'\circ\eta_7 .
\end{equation}

\paragraph*{v). The Group \(\pi^{n}_{n+5}\).}

\begin{proposition}\label{prop:5.9}
\ \\
% Prop5.9.\
\(\pi^{2}_{7}=\mathbb Z_2\gen{\eta_{2}\circ\nu'\circ\eta_6},\) \\
\(\pi^{3}_{8}=\mathbb Z_2\gen{\nu'\circ\eta_6^2},\) \\
\(\pi^{4}_{9}=\mathbb Z_2\gen{\nu_{4}\circ\eta_7^2}\oplus \mathbb Z_2\gen{E\nu'\circ\eta_7^2},\) \\
\(\pi^{5}_{10}=\mathbb Z_2\gen{\nu_{5}\circ\eta_8^2},\) \\
\(\pi^{6}_{11}=\mathbb Z\gen{\Delta(\iota_{13})},\) \\
\(\pi^{n}_{n+5}=0\ (n\ge7),\) \
\((G_{5};2)=0.\)
\end{proposition}

\noindent
\textbf{Proof.}\ 
(5.2) とProp5.8 から, 最初の結果を得る. \\
Lem5.7 により, \(E\pi_7^2=\gen{E(\eta_2\circ\nu')\circ\eta_7}=0 .\)
よって, (4.4) から, 
\[
0\to \pi^{3}_{8}\xrightarrow{H}\pi^{5}_{8}\xrightarrow{\Delta}\pi^{2}_{6}
\]
は完全列である. 
Lem5.7, Prop 5.6, (5.5) により, \(\ker\Delta\) は \(4\nu_5=\eta_5^3\) で生成される. \\
Prop2.2 と (5.3) により, 次を得る. 
\[
H(\nu'\circ \eta_6^2)=H(\nu')\circ \eta_6^2=\eta_5\circ \eta_6^2=\eta_5^{3}.
\]
よって, 完全性から, Propの2番目が成り立つ. \\
(5.6) により, \(\pi_9^4\) の結果を得る. 

\switchcolumn

\noindent
\textbf{(5.5).}\
\(n\ge5\ \text{に対して, } 2\nu_n=E^{\,n-3}\nu',\
4\nu_{n}=\eta_n^{3},\ 4\nu=\eta^{3}.\)

\noindent
\textbf{Prop3.1.}\
\(
\alpha \wedge \beta
= (-1)^{(p+k)h} E^{q}\alpha \circ E^{p+k}\beta
= (-1)^{ph} E^{p}\beta \circ E^{q+h}\alpha .
\)

\noindent
\textbf{(5.9).}\
\(\eta_3\circ \nu_4= \nu'\circ \eta_6,\
\eta_n\circ \nu_{n+1}=0\ (n\ge5),\
\nu_n\circ \eta_{n+3}=0\ (n\ge6)\)

\noindent
\textbf{(4.5).}\
\(E^{m-n}:\pit{n+k}{n}\xrightarrow{\cong}\pit{m+k}{m},\
(n\ge k+2). \)

\noindent
\textbf{Prop5.8.}\
\(\pi^{2}_{6}=\mathbb Z_4\gen{\eta_{2}\circ\nu'},\) \\
\(\pi^{3}_{7}=\mathbb Z_2\gen{\nu'\circ\eta_6},\) \\
\(\pi^{4}_{8}=\mathbb Z_2\gen{\nu_{4}\circ\eta_7}\oplus \mathbb Z_2\gen{E\nu'\circ\eta_7},\) \\
\(\pi^{5}_{9}=\mathbb Z_2\gen{\nu_{5}\circ\eta_8},\) \\
\(\pi^{n}_{n+4}=0\ (n\ge6),\) \
\((G_{4};2)=0.\)

\noindent
\textbf{(5.2).}\
\(\text{合成 } \alpha\longmapsto \eta_2\circ \alpha \text{ は, 同型 } 
\pi^{3}_{i}\xrightarrow{\cong}\pi^{2}_{i}\ (i\ge3) 
\text{ を定める. }\)

\noindent
\textbf{Lem5.7.}\
\(\alpha\in\pi_i(S^3)\) に対して,
\(E^2\alpha\in2\iota_5\circ\pi_{i+2}(S^5)\) ならば, 
\(E(\eta_2\circ\alpha)=0\) である. \\
特に, \(E(\eta_2\circ\nu')=0, \ \Delta(\nu_5)=\pm(\eta_2\circ\nu')\) である. 

\noindent
\textbf{Prop5.6.}\
\(\pi^{2}_{5}=\gen{\eta_{2}^{3}}\cong \mathbb Z_{2},\) \\
\(\pi^{3}_{6}=\gen{\nu'}\cong \mathbb Z_{4},\) \\
\(\pi^{4}_{7}=\gen{\nu_{4}}\oplus \gen{E\nu'}\cong \mathbb Z\oplus\mathbb Z_{4},\) \\
\(\pi^{n}_{n+3}=\gen{\nu_{n}}\cong \mathbb Z_{8}\ (n\ge5),\)\qquad
\((G_{3};2)=\gen{\nu}\cong \mathbb Z_{8}.\)

\noindent
\textbf{(5.5).}\
\(n\ge5\ \text{に対して, } 2\nu_n=E^{\,n-3}\nu',\
4\nu_{n}=\eta_n^{3},\ 4\nu=\eta^{3}.\)

\noindent
\textbf{Prop2.2.}\
$K,L$ を有限 cell complex とする. \\
$\alpha\in[EK, S^{m+1}], \ \beta\in[L,K], \ \gamma\in\pi_r(S^r)$ 
に対して, 次が成り立つ：\\
\(
H(\alpha\circ E\beta)\ =\ H(\alpha)\circ E\beta,
\)\\
\(
H(E\gamma\circ\alpha)\ =\ E(\gamma\wedge\gamma)\circ H(\alpha).
\)

\noindent
\textbf{(5.3).}\
\(\nu'\in\pi_6^3,\ H(\nu')=\eta_5,\ 2\nu'=\eta_3\circ\eta_4\circ\eta_5 .\)

\noindent
\textbf{(5.6).}\
\((\alpha,\beta)\ \longmapsto\ E\alpha+\nu_{4}\circ\beta
\quad\text{は,}\quad
\pi^{\,3}_{i-1}\oplus \pi^{\,7}_{i}\ \xrightarrow{\ \cong\ }\ \pi_{\,i}^{4}
\ \text{の同型を与える. }\)



\end{paracol}

\begin{paracol}{2}

(4.4) の次の完全列を考える. 
\[
\pi^{9}_{11}\xrightarrow{\ \Delta\ }\pi^{4}_{9}\xrightarrow{\ E }\pi^{5}_{10}
\xrightarrow{\ H }\pi^{9}_{10}\xrightarrow{\ \Delta\ }\pi^{4}_{8}\xrightarrow{E}\pi_9^5 .
\]
Prop5.8 と (5.11) により, \(\ker(E:\pi_8^4\to\pi_9^5)\cong\mathbb Z_2\) である. \\
\(\pi_{10}^9\cong\mathbb Z_2\) は, \(\Delta\) により, \(\pi_8^4\) に単射で写る. 
よって, \(E:\pi_9^4\to\pi_{10}^5\) は全射である. \\
\(\pi_{11}^9\) は \(\eta_9^2\) で生成され, Prop2.5 と (5.11) により, 
\(\Delta(\eta_9^2)=\Delta(\eta_9)\circ\eta_8=E\nu'\circ\eta_7^2\) である. 
よって, \(\pi_{10}^5=E\pi_9^4=\mathbb Z_2\gen{\nu_5\circ\eta_8^2}\) である. \\
(5.9) により, \(E(\nu_5\circ\eta_8^2)=\nu_6\circ\eta_9\circ\eta_{10}=0\) なので, 
\(E\pi_{10}^5=0 .\)\\
また, Prop2.5 と (5.10) により, \(\Delta(\eta_{11})=\nu_5\circ\eta_8^2\) である. 
よって, (4.4) から, 
\[
0\to \pi^{6}_{11}\xrightarrow{H}\pi^{11}_{11}\xrightarrow{\ \Delta\ }\pi^{5}_{9}
\]
が完全列である. 
Prop5.8 と (5.10) により, \(\Delta\) は全射で, \(\ker\Delta=\gen{2\iota_{11}}\) である. \\
Prop2.7 により, \(H(\Delta(\iota_{13}))=\pm2\iota_{11}\) なので, 
\(\pi_{11}^6=\mathbb Z \gen{\Delta(\iota_{13})}\) である. 
\[
\pi^{13}_{13}\xrightarrow[\cong]{\Delta}\pi^{6}_{11}\xrightarrow{E}\pi^{7}_{12}\xrightarrow{}\pi_{12}^{13}=0
\]
の完全性により, \(\pi_{12}^7=0\) である. 
(4.5) により, Propの最後の結果が得られる. 
\qed

証明の中で次が示された. 
\begin{equation}\label{eq:5.12}
% (5.12).\
\Delta:\ \pi^{2n+1}_{n+7}\ \longrightarrow\ \pi^{n}_{n+5}
\quad\text{は, }n=4,5,6\text{ で単射である. }
\end{equation}

次のLemを示す. 
\begin{lemma}\label{lemma:5.10}
% Lem5.10.\
\[
\Delta(\iota_{13})\in\toda{\nu_6,\eta_9,2\iota_{10}}\in\pi_{11}(S^6)/2\pi_{11}(S^6) .
\]
\end{lemma}

\noindent
\textbf{Proof.}\ 
この Toda bracket は, 次の群の coset である. (cf. (4.7), (5.9)).
\[
\nu_6\circ \pi_{11}(S^9)+\pi_{11}(S^6)\circ2\iota_{11}
= \nu_6\circ\pi_{11}^9+2\pi_{11}(S^6)=2\pi_{11}(S^6) .
\]
(5.10) とProp2.6 により, 
\(H\toda{\nu_6,\eta_9,2\iota_{10}}\ni\iota_{11}\circ E(2\iota_{10})=2\iota_{11} .\)\\
また, \(H(\Delta(\iota_{13}))=\pm2\iota_{11}\) なので, (2.11) の完全列
\[
\pi_{10}(S^5)\xrightarrow{E}\pi_{11}(S^6)\xrightarrow{H}\pi_{11}(S^{11})
\]
から, 
\(\Delta(\iota_{13})\in\toda{\nu_6,\eta_9,2\iota_{10}} + E\pi_{10}(S^5)\) である. \\
\(E\pi_{10}(S^5)\) は有限で, \(E\pi_{10}^5=0\) なので, 
\(E\pi_{10}(S^5)\subset2\pi_{11}(S^6) .\)\\
よって, \(\pm\Delta(\iota_{13})\in\toda{\nu_6,\eta_9,2\iota_{10}}\) である. 
\qed

\switchcolumn

\noindent
\textbf{Prop5.8.}\
\(\pi^{2}_{6}=\mathbb Z_4\gen{\eta_{2}\circ\nu'},\) \\
\(\pi^{3}_{7}=\mathbb Z_2\gen{\nu'\circ\eta_6},\) \\
\(\pi^{4}_{8}=\mathbb Z_2\gen{\nu_{4}\circ\eta_7}\oplus \mathbb Z_2\gen{E\nu'\circ\eta_7},\) \\
\(\pi^{5}_{9}=\mathbb Z_2\gen{\nu_{5}\circ\eta_8},\) \\
\(\pi^{n}_{n+4}=0\ (n\ge6),\) \
\((G_{4};2)=0.\)

\noindent
\textbf{(5.11).}\
\(\Delta(\eta_9)=E\nu'\circ\eta_7 .\)

\noindent
\textbf{Prop2.5.}\
$\alpha\in\pi_{i+1}(S^{2m+1}), \ \beta\in \pi_{j-1}(S^{i-1})$ 対して, 
\(
\Delta\alpha\ \circ\ \beta\ \subset\ \Delta(\alpha\circ E^2\beta).
\)\\
特に, 任意の $\gamma$ に対し,
\(
\Delta(E^2\gamma)\ \ni\ \pm\,[\,\iota_m,\ \iota_m\,]\circ\gamma
\)
が成り立つ. \\
ここで $[\, ,\,]$ は Whitehead product を表す. \\
上の「$\subset$」および「$\ni$」は, $m$ が奇数, 
あるいは $j<3m-1$, \\
または, 関係する元が \(2\)-primary 成分にあるときには, 
等号「$=$」に置き換えられる. 

\noindent
\textbf{(5.10).}\
\(\Delta(\iota_{11})=\nu_5\circ\eta_8 .\)

\noindent
\textbf{Prop2.7.}\
\(H\Delta(\iota_{2n+1}) = \pm2\iota_{2m-1},\ (n:\text{even})\).

\noindent
\textbf{(4.5).}\
\(E^{m-n}:\pit{n+k}{n}\xrightarrow{\cong}\pit{m+k}{m},\
(n\ge k+2). \)

\noindent
\textbf{(4.7).}\
\(n\ge1, \ \alpha\in\pis{k+n}{S^m; 2},\ \gamma\in\pis{i}{S^j; 2}
    \ \text{のとき, }\)\\
\(\alpha\circ E^n\pis{j}{S^k} = \alpha\circ E^n \pit{j}{k}
  \subset \pit{j+n}{m},
  \ \pis{j}{S^k}\circ\gamma = \pit{j}{k}\circ\gamma \subset\pit{i}{k}\ . \)\\
\(\text{さらに, }\ \alpha\circ E^n\beta=0,\ \beta\circ\gamma = 0 \ \text{のとき, }\)\\
\(\todai[n]{\alpha, E^n\beta, E^n\gamma}\ \text{は,}\
  \alpha\circ E^n\pit{i+1}{k} + \pit{j+n+1}{m}\circ E^{n+1}\gamma\
  \text{の coset で, }\pit{i+n+1}{m} \ \text{の部分集合である. }\)

\noindent
\textbf{(5.9).}\
\(\eta_3\circ \nu_4= \nu'\circ \eta_6,\
\eta_n\circ \nu_{n+1}=0\ (n\ge5),\
\nu_n\circ \eta_{n+3}=0\ (n\ge6)\)

\noindent
\textbf{Prop2.6.}\
\(H\todai[1]{E\alpha, E\beta, E\gamma} 
= -\Delta^{-1}(\alpha\circ\beta)\circ E^2\gamma.\)




\end{paracol}

\begin{paracol}{2}

% \bigskip
\paragraph*{vi). The Group \(\pi^{n}_{n+6}\).}
\[
\nu_n^2:=\nu_n\circ\nu_{n+3},(n\ge4)\qquad
\nu^2:=\nu\circ\nu.
\]

\begin{proposition}\label{prop:5.11}
\ \\
% Prop5.11.\
\(\pi^{2}_{8}=\mathbb Z_2\gen{\eta_{2}\circ\nu'\circ\eta_6^2},\) \\
\(\pi^{3}_{9}=0,\) \\
\(\pi^{4}_{10}=\mathbb Z_8\gen{\nu_{4}^2},\) \\
\(\pi^{n}_{n+6}=\mathbb Z_2\gen{\nu_n^2}\ (n\ge5),\) \
\((G_{6};2)=\mathbb Z_2\gen{\nu^2}.\)
\end{proposition}

\noindent
\textbf{Proof.}\ 
最初の結果は, Prop5.9 と (5.2) により証明される. \\
Lem5.7 により, \(E\pi_8^2=\gen{E(\eta_2\circ\nu')\circ\eta_6^2}=0 .\)\\
同様に, \(E\pi_7^2=0\) である. 
(4.4) により, 
\[
0\to \pi^{3}_{9}\xrightarrow{H}\pi^{5}_{9}\xrightarrow{\Delta}\pi^{2}_{7}\to 0
\]
は完全である. 
Prop \ref{prop:5.8}, \ref{prop:5.9} より,
\(\Delta\) は同型で, 
\(\pi^{3}_{9}=0 .\)\\
さらに, Prop \ref{prop:5.6} と (\ref{eq:5.6}) により,
\(\pi^{4}_{10}=\gen{\nu_4^{2}}\cong \mathbb Z_8 .\)\\
次の完全列を考える. 
\[
\pi^{2n+1}_{n+8}\xrightarrow{\Delta}\pi^{n}_{n+6}
\xrightarrow{E}\pi^{n+1}_{n+7}\xrightarrow{H}\pi^{2n+1}_{n+7}
\xrightarrow{P}\pi_{n+5}^n, 
\qquad(n=4,5,6) .
\]
(5.12) により, \(H\) は自明, \(E\) は全射である. 
ここで, 次を証明する. 
\begin{equation}\label{eq:5.13}
% (5.13).\
\Delta(\nu_9)=\pm\,2(\nu_4^{2}),\qquad
\Delta(\eta^{2}_{11})=0,\qquad
\Delta(\eta_{13})=0.
\end{equation}

Prop 2.5 と (\ref{eq:5.8}) より,
\[
\Delta(\nu_9)=\Delta(\iota_9)\circ \nu_7=\pm\,(2\nu_4-E\nu')\circ \nu_7
=\pm\big(2\,\nu_4^{2}-E(\nu'\circ \nu_6)\big).
\]
\(\nu'\circ \nu_6\in \pi^{3}_{9}=0\) なので, 
\(\Delta(\nu_9)=\pm 2(\nu_4^{2}) .\) また,
\[
\begin{align*}
\Delta(\eta^{2}_{11})
& =  \Delta(\iota_{11})\circ \eta_9^2
&& \text{by Prop2.5}\\
& =  \nu_5\circ \eta_8\circ \eta_9^2
= \nu_5\circ \eta_8^{3}
&& \text{by (5.10)}\\
& =  4(\nu_5\circ\nu_8)
= 4E(\nu_4^2)
&& \text{by (5.5)}\\
& =  E\Delta(\pm2\nu_9)
=0 .
\end{align*}
\]
\[
\begin{align*}
\Delta(\eta_{13})
& =  \Delta(\iota_{13})\circ \eta_{11}
&& \text{by Prop2.5}\\
& \in \toda{ \nu_6,\ \eta_9,\ 2\iota_{10} }\circ \eta_{11}
&& \text{by Lem5.10}\\
& =  \nu_6\circ\toda{\eta_9,2\iota_{10},\eta_{10}} 
&& \text{by Prop1.4}.
\end{align*}
\]

\switchcolumn

\noindent
\textbf{Prop5.9.}\
\(\pi^{2}_{7}=\mathbb Z_2\gen{\eta_{2}\circ\nu'\circ\eta_6},\) \\
\(\pi^{3}_{8}=\mathbb Z_2\gen{\nu'\circ\eta_6^2},\) \\
\(\pi^{4}_{9}=\mathbb Z_2\gen{\nu_{4}\circ\eta_7^2}\oplus \mathbb Z_2\gen{E\nu'\circ\eta_7^2},\) \\
\(\pi^{5}_{10}=\mathbb Z_2\gen{\nu_{5}\circ\eta_8^2},\) \\
\(\pi^{6}_{11}=\mathbb Z_2\gen{\Delta(\iota_{13})},\) \\
\(\pi^{n}_{n+5}=0\ (n\ge7),\) \
\((G_{5};2)=0.\)

\noindent
\textbf{(5.2).}\
\(\text{合成 } \alpha\longmapsto \eta_2\circ \alpha \text{ は, 同型 } 
\pi^{3}_{i}\xrightarrow{\cong}\pi^{2}_{i}\ (i\ge3) 
\text{ を定める. }\)

\noindent
\textbf{Lem5.7.}\
\(\alpha\in\pi_i(S^3)\) に対して,
\(E^2\alpha\in2\iota_5\circ\pi_{i+2}(S^5)\) ならば, 
\(E(\eta_2\circ\alpha)=0\) である. \\
特に, \(E(\eta_2\circ\nu')=0, \ \Delta(\nu_5)=\pm(\eta_2\circ\nu')\) である. 

\noindent
\textbf{Prop5.8.}\
\(\pi^{2}_{6}=\mathbb Z_4\gen{\eta_{2}\circ\nu'},\) \\
\(\pi^{3}_{7}=\mathbb Z_2\gen{\nu'\circ\eta_6},\) \\
\(\pi^{4}_{8}=\mathbb Z_2\gen{\nu_{4}\circ\eta_7}\oplus \mathbb Z_2\gen{E\nu'\circ\eta_7},\) \\
\(\pi^{5}_{9}=\mathbb Z_2\gen{\nu_{5}\circ\eta_8},\) \\
\(\pi^{n}_{n+4}=0\ (n\ge6),\) \
\((G_{4};2)=0.\)

\noindent
\textbf{Prop5.6.}\
\(\pi^{2}_{5}=\gen{\eta_{2}^{3}}\cong \mathbb Z_{2},\) \\
\(\pi^{3}_{6}=\gen{\nu'}\cong \mathbb Z_{4},\) \\
\(\pi^{4}_{7}=\gen{\nu_{4}}\oplus \gen{E\nu'}\cong \mathbb Z\oplus\mathbb Z_{4},\) \\
\(\pi^{n}_{n+3}=\gen{\nu_{n}}\cong \mathbb Z_{8}\ (n\ge5),\)\qquad
\((G_{3};2)=\gen{\nu}\cong \mathbb Z_{8}.\)

\noindent
\textbf{(5.6).}\
\((\alpha,\beta)\ \longmapsto\ E\alpha+\nu_{4}\circ\beta
\quad\text{は,}\quad
\pi^{\,3}_{i-1}\oplus \pi^{\,7}_{i}\ \xrightarrow{\ \cong\ }\ \pi_{\,i}^{4}
\ \text{の同型を与える. }\)

\noindent
\textbf{(5.12).}\
\(\Delta:\ \pi^{2n+1}_{n+7}\ \longrightarrow\ \pi^{n}_{n+5}
\quad\text{は, }n=4,5,6\text{ で単射である. }\)

\noindent
\textbf{Prop2.5.}\
$\alpha\in\pi_{i+1}(S^{2m+1}), \ \beta\in \pi_{j-1}(S^{i-1})$ 対して, 
\(
\Delta\alpha\ \circ\ \beta\ \subset\ \Delta(\alpha\circ E^2\beta).
\)\\
特に, 任意の $\gamma$ に対し,
\(
\Delta(E^2\gamma)\ \ni\ \pm\,[\,\iota_m,\ \iota_m\,]\circ\gamma
\)
が成り立つ. \\
ここで $[\, ,\,]$ は Whitehead product を表す. \\
上の「$\subset$」および「$\ni$」は, $m$ が奇数, 
あるいは $j<3m-1$, \\
または, 関係する元が \(2\)-primary 成分にあるときには, 
等号「$=$」に置き換えられる. 

\noindent
\textbf{(5.8).}\
\(\Delta(\iota_9)=\pm(2\nu_4-E\nu')  \ (=\pm[\iota_4, \iota_4]) .\)

\noindent
\textbf{(5.10).}\
\(\Delta(\iota_{11})=\nu_5\circ\eta_8 .\)

\noindent
\textbf{(5.5).}\
\(n\ge5\ \text{に対して, } 2\nu_n=E^{\,n-3}\nu',\
4\nu_{n}=\eta_n^{3},\quad 4\nu=\eta^{3}.\)

\noindent
\textbf{Lem5.10.}\
\(
\Delta(\iota_{13})\in\toda{\nu_6,\eta_9,2\iota_{10}}\in\pi_{11}(S^6)/2\pi_{11}(S^6) .
\)

\noindent
\textbf{Prop1.4.}\
\(\todai[n]{\alpha, E^n\beta, E^n\gamma}\circ E^{n+1}\delta
= (-1)^{n+1}(\alpha\circ E^n\toda{\beta, \gamma, \delta})\).



\end{paracol}

\begin{paracol}{2}

(5.4) により, 
\(\nu_6\circ\gen{\eta_9,2\iota_{10},\eta_{10}}
=\gen{\nu_6\circ E^6\nu', -\nu_6\circ E^6\nu'} .\)\\
(5.5) より, 
\(\nu_6 E^6\nu'=2\nu_6^2=E^2\Delta(\nu_9)=0\)
なので, \(\Delta(\eta_{13})=0\) となり, (5.13) が示された. \\
(5.13) を, 上の完全列に適用することにより, \(n=5,6,7\) に対してPropは真である. \\
Prop4.4 により, 
\(\pi_{14}^8\) は 
\(\pi_{13}^7\oplus\pi_{14}^{15}=\pi_{13}^7\)
に, \(E\) で同型となる. \\
これで, Propの \(n=8\) の場合が示された. \\
\(n\geq9\) に対しては, stable の範囲なので, (4.5) により, 
Prop5.11 が証明できた. 
\qed

\begin{lemma}\label{lemma:5.12}
% \noindent
% \textbf{Lem5.12.}\
\(\toda{\eta_n,\nu_{n+1},\eta_{n+4}} = \gen{\nu_n^2},\ (n\geq6) .\)
\end{lemma}

\noindent
\textbf{Proof.}\ 
(5.9) により, この Toda bracket は定義され, \\
\(\eta_n\circ \pi_{n+6}(S^{n+1})+\pi_{n+5}(S^{n})\circ \eta_{n+5}\) 
の coset である. \\
(3.2) により, \(\pi_{n+6}(S^{n+1})=E\pi_{n+5}(S^{n})\), \\
(4.7) とProp\ref{prop:5.9} から, 
\(\eta_n\circ \pi_{n+6}(S^{n+1})=\eta_n\circ \pi^{n}_{n+6}=0 .\)\\
同様に,
\(\pis{n+5}{S^n}\circ \eta_{n+5}=\pi_{n+5}^n\circ\eta_{n+5}\)
であり, \\
この群は \(n>6\) で自明で, 
\(n=6\) のとき, \(\Delta(\iota_{13})\circ\eta_{11}\) で生成される. \\
(5.13) の証明の中で見たように, \(\Delta(\iota_{13})\circ\eta_{11}=\Delta(\eta_{13})=0 .\)\\
全ての場合で, この Toda bracket は, 単一の元
\(x\,\nu_n^{2}\) (\(x=0\) または \(1\))から成る. \\
(1.15) とProp 1.3により, 
\(
E\toda{\eta_n,\nu_{n+1},\eta_{n+4}}
\subset -\toda{\eta_{n+1},\nu_{n+2},\eta_{n+5}} .
\)\\
これは, \(x\) の値が \(n\) によらないことを示している. \\
ここで, \(\beta=\nu_m \ (m\geq6)\) に対して, Lem\ref{lem:5.5} を適用する. \\
\(
E^{2}\nu_m\circ E^{m+1}\nu_4
=\nu_{m+2}^{2}
\in \todai[3]{\eta_{m+2},E^3\nu_{m},\eta_{m+6}}
\subset \toda{\eta_{m+2},\nu_{m+3},\eta_{m+6}} .
\)\\
よって, \(\nu_{m+2}^2=x\nu_{m+2}^2,\ x=1\) である. 
\qed

Remark.\  Lem5.12 は, \(n=5\) の場合も成立する. 


\paragraph*{vii). The Group \(\pi^{n}_{n+7}\).}

(\ref{eq:5.2}) とProp \ref{prop:5.11} より,
\(\pi^{2}_{9}=0\). 
(4.4) から, 次の列は完全である. 
\[
0=\pi^{2}_{9}\xto
\pi^{3}_{10}\xrightarrow{H}
\pi_{10}^5\xrightarrow{\Delta}
\pi^{2}_{8}\xrightarrow{E}\pi^{3}_{9}=0 .
\]
Prop \ref{prop:5.9} とProp \ref{prop:5.11} より,
\(\Delta\) は同型で, 
\(\pi^{3}_{10}=0\). \\
さらに (\ref{eq:5.6}) とProp \ref{prop:5.8} から, 
\(\pi^{4}_{11}\cong \pi^{3}_{10}\oplus \pi^{7}_{11}=0\). \\
ここで, 次の完全列を考える. 
\[
0=\pi^{4}_{11}\xto\pi^{5}_{12}\xrightarrow{H}
\pi^{9}_{12}\xrightarrow{\Delta}\pi^{4}_{10} .
\]
(\ref{eq:5.13}), Prop \ref{prop:5.11}, Prop \ref{prop:5.6} から, 
\(\Delta\) は 次数 2 の準同型 で, \\
その kernel は \(4\nu_9=\eta_9^{3}\) により生成される.
よって, 同型 
\(H:\pi^{5}_{12}\xrightarrow{}4\pi^{9}_{12}\ \cong\ \mathbb Z_2 .\)
を得る. 

\switchcolumn

\noindent
\textbf{(5.4).}\
\(\todai[t]{\eta_{n},\,2\iota_{n+1},\,\eta_{n+1}} = \{E^{\,n-3}\nu',-E^{\,n-3}\nu'\}.\)

\noindent
\textbf{(5.5).}\
\(n\ge5\ \text{に対して, } 2\nu_n=E^{\,n-3}\nu',\
4\nu_{n}=\eta_n^{3},\quad 4\nu=\eta^{3}.\)

\noindent
\textbf{(5.13).}\
\(\Delta(\nu_9)=\pm\,2(\nu_4^{2}),\
\Delta(\eta^{2}_{11})=0,\
\Delta(\eta_{13})=0.\)

\noindent
\textbf{Prop5.11.}\
\(\pi^{2}_{8}=\mathbb Z_2\gen{\eta_{2}\circ\nu'\circ\eta_6^2},\)\quad 
\(\pi^{3}_{9}=0,\)\quad \(\pi^{4}_{10}=\mathbb Z_8\gen{\nu_{4}^2},\) \\
\(\pi^{n}_{n+6}=\mathbb Z_2\gen{\nu_n^2}\ (n\ge5),\) \
\((G_{6};2)=\mathbb Z_2\gen{\nu^2}.\)

\noindent
\textbf{Prop4.4.}\
\(H(\alpha) = \pm 2\iota_{2m-1}\ \Longrightarrow\
\pit{i-1}{m-1}\oplus\pit{i}{2m-1}\xrightarrow{\cong}\pit{i}{m},\
(\beta,\gamma)\mapsto E\beta+\alpha\circ\gamma.\)\\
\(0=\pit{2}{1}\to\pit{3}{2}
\xrightarrow{H}{\cong}\pit{3}{3}
\xrightarrow{\Delta}\pit{1}{1}
\xrightarrow{E}\pit{2}{2}.\)

\noindent
\textbf{(4.5).}\
\(E^{m-n}:\pit{n+k}{n}\xrightarrow{\cong}\pit{m+k}{m},\
(n\ge k+2). \)

\noindent
\textbf{(5.9).}\
\(\eta_3\circ \nu_4= \nu'\circ \eta_6,\
\eta_n\circ \nu_{n+1}=0\ (n\ge5),\
\nu_n\circ \eta_{n+3}=0\ (n\ge6)\)

\noindent
\textbf{(3.2).}\
\(\Ex: \pis{i}{S^m} \longrightarrow \pis{i+1}{S^{m+1}}
\quad\text{は, } i<2m-1 \text{で同型, } i=2m-1 \text{で全射である. }\)

\noindent
\textbf{(4.7).}\
\(n\ge1, \ \alpha\in\pis{k+n}{S^m; 2},\ \gamma\in\pis{i}{S^j; 2}
    \ \text{のとき, }\)\\
\(\alpha\circ E^n\pis{j}{S^k} = \alpha\circ E^n \pit{j}{k}
  \subset \pit{j+n}{m},
  \ \pis{j}{S^k}\circ\gamma = \pit{j}{k}\circ\gamma \subset\pit{i}{k}\ . \)\\
\(\text{さらに, }\ \alpha\circ E^n\beta=0,\ \beta\circ\gamma = 0 \ \text{のとき, }\)\\
\(\todai[n]{\alpha, E^n\beta, E^n\gamma}\ \text{は,}\
  \alpha\circ E^n\pit{i+1}{k} + \pit{j+n+1}{m}\circ E^{n+1}\gamma\
  \text{の coset で, }\pit{i+n+1}{m} \ \text{の部分集合である. }\)

\noindent
\textbf{Prop5.9.}\
\(\pi^{2}_{7}=\mathbb Z_2\gen{\eta_{2}\circ\nu'\circ\eta_6},\) \quad
\(\pi^{3}_{8}=\mathbb Z_2\gen{\nu'\circ\eta_6^2},\) \quad
\(\pi^{4}_{9}=\mathbb Z_2\gen{\nu_{4}\circ\eta_7^2}\oplus \mathbb Z_2\gen{E\nu'\circ\eta_7^2},\) \\
\(\pi^{5}_{10}=\mathbb Z_2\gen{\nu_{5}\circ\eta_8^2},\) \quad
\(\pi^{6}_{11}=\mathbb Z_2\gen{\Delta(\iota_{13})},\) \quad
\(\pi^{n}_{n+5}=0\ (n\ge7),\) \
\((G_{5};2)=0.\)

\noindent
\textbf{(1.15).}\
\(\{\alpha,\ \Ex^n\beta,\ \Ex^n\gamma\}_n
\ \subset\
\bigl\{\alpha,\ \Ex^{m}(E^{n-m}\beta),\ E^{m}(E^{n-m}\gamma)\bigr\}_m .\)

\noindent
\textbf{Prop1.3.}\
\(-E\todai[n]{\alpha, E^n\beta, E^n\gamma}
\subset \todai[n+1]{E\alpha, E^{n+1}\beta, E^{n+1}\gamma}\).

\noindent
\textbf{Lem5.5.}\
$Lem 5.4$ の元 \(\nu_{4}\) に対し, 
Toda bracket 
\(\todai[3]{\eta_{m+2},\,E^{3}\beta,\,\eta_{t+5}}\) は,\\
\(
E^{2}\beta\circ E^{t}\nu_{4}\text{または,} -E^{2}\beta\circ E^{t}\nu_{4}
\)
を含む. \\
ここで, \(\beta\in\pi_{t+2}(S^{m})\) は,
\(\beta\circ \eta_{t+2}=0\), \(t>0\) を満たす元である. 

\noindent
\textbf{(5.2).}\
\(\text{合成 } \alpha\longmapsto \eta_2\circ \alpha \text{ は, 同型 } 
\pi^{3}_{i}\xrightarrow{\cong}\pi^{2}_{i}\ (i\ge3) 
\text{ を定める. }\)

\noindent
\textbf{(5.6).}\
\((\alpha,\beta)\ \longmapsto\ E\alpha+\nu_{4}\circ\beta
\quad\text{は,}\quad
\pi^{\,3}_{i-1}\oplus \pi^{\,7}_{i}\ \xrightarrow{\ \cong\ }\ \pi_{\,i}^{4}
\ \text{の同型を与える. }\)

\noindent
\textbf{Prop5.8.}\
\(\pi^{2}_{6}=\mathbb Z_4\gen{\eta_{2}\circ\nu'},\) \quad
\(\pi^{3}_{7}=\mathbb Z_2\gen{\nu'\circ\eta_6},\) \quad
\(\pi^{4}_{8}=\mathbb Z_2\gen{\nu_{4}\circ\eta_7}\oplus \mathbb Z_2\gen{E\nu'\circ\eta_7},\) \\
\(\pi^{5}_{9}=\mathbb Z_2\gen{\nu_{5}\circ\eta_8},\) \quad
\(\pi^{n}_{n+4}=0\ (n\ge6),\) \
\((G_{4};2)=0.\)

\noindent
\textbf{Prop5.6.}\
\(\pi^{2}_{5}=\gen{\eta_{2}^{3}}\cong \mathbb Z_{2},\) \quad
\(\pi^{3}_{6}=\gen{\nu'}\cong \mathbb Z_{4},\) \quad
\(\pi^{4}_{7}=\gen{\nu_{4}}\oplus \gen{E\nu'}\cong \mathbb Z\oplus\mathbb Z_{4},\) \\
\(\pi^{n}_{n+3}=\gen{\nu_{n}}\cong \mathbb Z_{8}\ (n\ge5),\)\quad
\((G_{3};2)=\gen{\nu}\cong \mathbb Z_{8}.\)



\end{paracol}

\begin{paracol}{2}

\begin{lemma}\label{lemma:5.13}
% Lem5.13.\
\(\todai[t]{\nu_5,8\iota_8,\nu_8},\ (0\leq t\leq3)\) 
は, ただ$1$つの元から成る. \\
この元を \(\sigma'''\) と書くと, 
\(H(\sigma''')=4\nu_9=\eta_9^3\) であり, \\
\(\pi_{12}^5=\mathbb{Z}_2\gen{\sigma'''}\)
である. 
\end{lemma}

\noindent
\textbf{Proof.}\ 
この Toda bracket は,
\(\nu_5\circ \pi_{12}(S^{8})+\pi_{9}(S^{5})\circ \nu_9\) の coset である. \\
(3.2),(4.7), Prop \ref{prop:5.8} から 
\[
\nu_5\circ\pis{12}{S^{8}}
=\nu_5\circ E\pis{11}{S^{7}}
=\nu_5\circ \pit{11}{7}=0.
\]
(3.2), Prop5.8, (5.9) により, 
\[
\pi_{9}(S^{5})\circ\nu_9=\pi_9^5\circ\nu_9
=\gen{\nu_5\circ \eta_8\circ \nu_9}=0 .
\]
よって, Toda bracket は単一元から成る. \\
さらに, (1.15) より,
\(\todai[t]{\nu_5,8\iota_8,\nu_8}
\subset \toda{ \nu_5,8\iota_8,\nu_8}\) である. \\
これから, \(\todai[t]{\nu_5,8\iota_8,\nu_8}\) も, 
1つの元 \(\sigma'''\) だけから成り, 
(4.7) により, \(\sigma'''\in\pi_{12}^5\) である. \\
Prop 2.6 から, 
\[
H(\sigma''')
= H\todai[1]{\nu_5, 8\iota_8, \nu_8}
= \Delta^{-1}(\nu_4\circ 8\iota_7)\circ \nu_9 .
\]
(5.8) により, 
\[
\Delta(4\iota_9)=\pm4(2\nu_4-E\nu')
=\pm(8\nu_4-E4\nu')
=\pm8\nu_4
=\pm(\nu_4\circ8\iota_7) .
\]
よって, 
\[
\Delta^{-1}(\nu_4\circ8\iota_7)\circ\nu_9
\ni\pm(4\iota_9)\circ\nu_9
=4\nu_9
=\eta_9^3 .
\]
従って, \(H(\sigma''')=4\nu_9\) であり, \\
同型 \(H:\pi^{5}_{12}
\xrightarrow{\cong}4\pi^{9}_{12}
\cong\mathbb Z_2\) により,
Lemが全て証明される. 
\qed

次の列は完全列である.
\begin{equation}\label{eq:5.14}
% (5.14).\
0\to \pi^{5}_{12}\xrightarrow{E}\pi^{6}_{13}\xrightarrow{H}\pi^{11}_{13}\to0,
\qquad
0\to \pi^{6}_{13}\xrightarrow{E}\pi^{7}_{14}\xrightarrow{H}\pi^{13}_{14}\to0.
\end{equation}

\noindent
これを示すには, 
\(\Delta(\pi_{14}^{11})=\Delta(\pi_{13}^{11})
=\Delta(\pi_{15}^{13})=\Delta(\pi_{14}^{13})=0\)
を証明すれば十分である. \\
(5.13) により, \(\Delta(\pi_{13}^{11})=\Delta(\pi_{14}^{13})=0\) である. \\
Prop2.5 と (5.13) により, 
\(\Delta(\eta_{13}^2)=\Delta(\eta_{13})\circ\eta_{12}=0 .\)\\
よって, \(\Delta(\pi_{15}^{13})=0\) である. \\
また, Prop2.5, (5.10), (5.9) により, 
\(\Delta(\nu_{11})=\Delta(\iota_{11})\circ\nu_9
=\nu_5\circ\eta_8\circ\nu_9=0 .\)\\
Prop5.6 により, \(\Delta(\pi_{14}^{11})=0 .\)\\
以上で (5.14) の完全性が示された. 

\switchcolumn

\noindent
\textbf{(3.2).}\
\(\Ex: \pis{i}{S^m} \longrightarrow \pis{i+1}{S^{m+1}}
\quad\text{は, } i<2m-1 \text{で同型, } i=2m-1 \text{で全射である. }\)

\noindent
\textbf{(4.7).}\
\(n\ge1, \ \alpha\in\pis{k+n}{S^m; 2},\ \gamma\in\pis{i}{S^j; 2}
    \ \text{のとき, }\)\\
\(\alpha\circ E^n\pis{j}{S^k} = \alpha\circ E^n \pit{j}{k}
  \subset \pit{j+n}{m},
  \ \pis{j}{S^k}\circ\gamma = \pit{j}{k}\circ\gamma \subset\pit{i}{k}\ . \)\\
\(\text{さらに, }\ \alpha\circ E^n\beta=0,\ \beta\circ\gamma = 0 \ \text{のとき, }\)\\
\(\todai[n]{\alpha, E^n\beta, E^n\gamma}\ \text{は,}\
  \alpha\circ E^n\pit{i+1}{k} + \pit{j+n+1}{m}\circ E^{n+1}\gamma\
  \text{の coset で, }\pit{i+n+1}{m} \ \text{の部分集合である. }\)

\noindent
\textbf{Prop5.8.}\
\(\pi^{2}_{6}=\mathbb Z_4\gen{\eta_{2}\circ\nu'},\) \quad
\(\pi^{3}_{7}=\mathbb Z_2\gen{\nu'\circ\eta_6},\) \quad
\(\pi^{4}_{8}=\mathbb Z_2\gen{\nu_{4}\circ\eta_7}\oplus \mathbb Z_2\gen{E\nu'\circ\eta_7},\) \\
\(\pi^{5}_{9}=\mathbb Z_2\gen{\nu_{5}\circ\eta_8},\) \quad
\(\pi^{n}_{n+4}=0\ (n\ge6),\) \
\((G_{4};2)=0.\)

\noindent
\textbf{(5.9).}\
\(\eta_3\circ \nu_4= \nu'\circ \eta_6,\
\eta_n\circ \nu_{n+1}=0\ (n\ge5),\
\nu_n\circ \eta_{n+3}=0\ (n\ge6)\)

\noindent
\textbf{(1.15).}\
\(\{\alpha,\ \Ex^n\beta,\ \Ex^n\gamma\}_n
\ \subset\
\bigl\{\alpha,\ \Ex^{m}(E^{n-m}\beta),\ E^{m}(E^{n-m}\gamma)\bigr\}_m .\)

\noindent
\textbf{Prop2.6.}\
\(H\todai[1]{E\alpha, E\beta, E\gamma} 
= -\Delta^{-1}(\alpha\circ\beta)\circ E^2\gamma.\)

\noindent
\textbf{(5.8).}\
\(\Delta(\iota_9)=\pm(2\nu_4-E\nu')  \ (=\pm[\iota_4, \iota_4]) .\)

\noindent
\textbf{(5.13).}\
\(\Delta(\nu_9)=\pm\,2(\nu_4^{2}),\
\Delta(\eta^{2}_{11})=0,\
\Delta(\eta_{13})=0.\)

\noindent
\textbf{Prop2.5.}\
$\alpha\in\pi_{i+1}(S^{2m+1}), \ \beta\in \pi_{j-1}(S^{i-1})$ 対して, 
\(
\Delta\alpha\ \circ\ \beta\ \subset\ \Delta(\alpha\circ E^2\beta).
\)\\
特に, 任意の $\gamma$ に対し,
\(
\Delta(E^2\gamma)\ \ni\ \pm\,[\,\iota_m,\ \iota_m\,]\circ\gamma
\)
が成り立つ. \\
ここで $[\, ,\,]$ は Whitehead product を表す. \\
上の「$\subset$」および「$\ni$」は, $m$ が奇数, 
あるいは $j<3m-1$, \\
または, 関係する元が \(2\)-primary 成分にあるときには, 
等号「$=$」に置き換えられる. 

\noindent
\textbf{(5.10).}\
\(\Delta(\iota_{11})=\nu_5\circ\eta_8 .\)

\noindent
\textbf{Prop5.6.}\
\(\pi^{2}_{5}=\gen{\eta_{2}^{3}}\cong \mathbb Z_{2},\) \quad
\(\pi^{3}_{6}=\gen{\nu'}\cong \mathbb Z_{4},\) \quad
\(\pi^{4}_{7}=\gen{\nu_{4}}\oplus \gen{E\nu'}\cong \mathbb Z\oplus\mathbb Z_{4},\) \\
\(\pi^{n}_{n+3}=\gen{\nu_{n}}\cong \mathbb Z_{8}\ (n\ge5),\)\quad
\((G_{3};2)=\gen{\nu}\cong \mathbb Z_{8}.\)




\end{paracol}

\begin{paracol}{2}

\begin{lemma}\label{lemma:5.14}
% Lem5.14.\
\(\sigma_{8}\in\pi^{8}_{15}\ , \ \sigma'\in\pi^{7}_{14}\ , \ \sigma''\in\pi^{6}_{13}\) 
が存在して, 次を満たす. \\
\(H(\sigma_{8})=\iota_{15} \ ,\qquad 2E\sigma_{8}=E^{2}\sigma',\)\\
\(H(\sigma')=\eta_{13} \ ,\qquad 2\sigma'=E\sigma'',\)\\
\(H(\sigma'')=\eta_{11}^2 \ ,\qquad 2\sigma''=E\sigma''' .\)
\end{lemma}

\noindent
\textbf{Proof.}\ 
Th 3.6 を \(\alpha=\nu_{4}\) に適用し, \\
Prop \ref{prop:5.11} により, 
\(\nu_4\) は (3.1) の条件を満たす. 
\[
2E^{4-3}\nu_{4}\circ E^{7-3}\nu_{4}=2\nu_{5}^{2}=0,
\quad (h=3).
\]
このとき, 
\(\alpha^{\ast}\in\pi^{8}_{15}\) が存在し, 
\(\beta\circ\nu_{t+4}=0\) を満たす任意の \(\beta\in\pi_{t+4}(S^m)\) に対して, 
\[
E^{4}\beta\circ E^t\alpha^{\ast}
\in (-1)^{m}\todai[7]{\nu_{m+4},\,E^{7}\beta,\,\nu_{t+11}}
+(-1)^{t}\todai[t+3]{E^{4}\beta,\,\nu_{t+8},\,2\nu_{t+11}} .
\]
今, \(\beta=8\iota_{5}\) とする. 
\(8\iota_5\circ\nu_5=8\nu_5=0\) なので, 
\[
8E\alpha^{\ast}=E^48\iota_5\circ E\alpha^{\ast}
\in -\todai[7]{\nu_9,8\iota_{12},\nu_{12}}
-\todai[4]{8\iota_9,\nu_9,2\nu_{12}} .  
\]
Toda bracket \(\todai[7]{\nu_{9},\,8\iota_{12},\,\nu_{12}}\) は, 
(4.7) とProp5.8 により, 
\[
\nu_9\circ E^7\pi_9(S^5) + \pi_{13}^9\circ\nu_{13}
=\nu_9\circ E^7\pi_9^5 + \pi_{13}^9\circ\nu_{13} = 0
\]
の coset なので, ただ1つの元である. 
Prop 1.3, Lem 5.13 により, 
\[
E\sigma'''\in E^4\todai[3]{\nu_5,8\iota_8,\nu_8}
\subset \todai[7]{\nu_9,8\iota_{12},\nu_{12}} .
\]
したがって, \(-\todai[7]{\nu_9,8\iota_{12},\nu_{12}}
=-E^4\sigma'''=E^4\sigma'''\) である. 
Prop1.2 により, 
\begin{align*}
\todai[4]{8\iota_9,\nu_9,2\nu_{12}}
& \subset\todai[4]{2\iota_9,4\iota_9\circ\nu_9\circ2\iota_{12}, \nu_{12}} \\
& =\todai[4]{2\iota_9,0,\nu_{12}}
=2\iota_9\circ E^4\pi_{12}(S^5) + \pi_{13}(S^9)\circ\nu_{13} \\
& =2E^4\pi_{12}(S^5) + \pi_{13}^9\circ\nu_{13}
=2E^4\pi_{12}(S^5) .
\end{align*}
以上から, \(8E\alpha^{\ast} - E^4\sigma''' \in 2E^4\pi_{12}(S^5)\) である. \\
\(\pi_{12}^5=\pi_{12}(S^5;2)\cong\mathbb{Z}_2\) なので, 
\(2E^4\pi_{12}(S^5)\) は奇数位数の有限群である. 
\[
8x E\alpha^{\ast}=xE^{4}\sigma'''=E^4\sigma'''\ , \ (x:\text{odd}) .
\]
(5.14) とLem5.13 により, \(\pi_{13}^6\) と \(\pi_{14}^7\) 
の位数は, それぞれ 4 と 8 であり, \\
\(E^2:\pi_{12}^5\to\pi_{14}^7\) は単射である. \\
また, (4.8) から, \(\pi_{16}^9\) の位数は 16 で, 
\(E^4:\pi_{12}^5\to\pi_{16}^9\) は単射である. \\
これから, \(8xE\alpha^{\ast}=E^4\sigma'''\ne0\) であり, 
\(16xE\alpha^{\ast}=0\) である. \\
これは, \(xE\alpha^{\ast}\) の位数が 16 で, 
\(\pi_{16}^9\cong\mathbb{Z}_{16}\) を生成することを示している. 

\switchcolumn

\noindent
\textbf{Th3.6.}\
$\alpha\in\pis{n+k}{S^{n}}$, $h\ge0$ が (3.1) を満たすとする.\\ 
さらに, $k\le 2n-2$ と仮定する.
このとき,
\(\alpha^{*}\in\pis{2n+2k+1}{S^{2n}}\)
が存在して, \\
\(\beta\circ \Ex^{\ell}\alpha=0\)
を満たす任意の $\beta\in\pis{n+\ell}{S^{n}}$ に対して, \\
\(\Ex^{n}\beta \circ \Ex^{t}\alpha^{*}
\in (-1)^{km+kt+t}\,\todai[n+k]{\Ex^{m}\alpha,\ 
  \Ex^{n+k}\beta,\ \Ex^{n+k+t}\alpha}\)\\
\(+ (-1)^{kn+h+t+1}\,\todai[h+t]{\Ex^{n}\beta,\ 
  \Ex^{n+t}\alpha,\ (1-(-1)^k)\Ex^{n+k+t}\alpha}\ .\)\\
さらに,
\((1-(-1)^k)\Ex^{n+h+1}\alpha=0\), 特に \(k\) が偶数のとき, \\
\(
\Ex^{n}\beta\circ \Ex^{t}\alpha^{*}\ \in\
(-1)^{km+kt+t}\,\todai[n+k]{\Ex^{m}\alpha,\ \Ex^{n+k}\beta,\ \Ex^{n+k+t}\alpha}\ .
\)

\noindent
\textbf{Prop5.11.}\
\(\pi^{2}_{8}=\mathbb Z_2\gen{\eta_{2}\circ\nu'\circ\eta_6^2},\)\quad 
\(\pi^{3}_{9}=0,\)\quad \(\pi^{4}_{10}=\mathbb Z_8\gen{\nu_{4}^2},\) \\
\(\pi^{n}_{n+6}=\mathbb Z_2\gen{\nu_n^2}\ (n\ge5),\) \
\((G_{6};2)=\mathbb Z_2\gen{\nu^2}.\)

\noindent
\textbf{(3.1).}\
\(\bigl(1-(-1)^k\bigr)\, \Ex^{\,n-h}\alpha \circ \Ex^{\,n+k-h}\alpha = 0.\)

\noindent
\textbf{(4.7).}\
\(n\ge1, \ \alpha\in\pis{k+n}{S^m; 2},\ \gamma\in\pis{i}{S^j; 2}
    \ \text{のとき, }\)\\
\(\alpha\circ E^n\pis{j}{S^k} = \alpha\circ E^n \pit{j}{k}
  \subset \pit{j+n}{m},
  \ \pis{j}{S^k}\circ\gamma = \pit{j}{k}\circ\gamma \subset\pit{i}{k}\ . \)\\
\(\text{さらに, }\ \alpha\circ E^n\beta=0,\ \beta\circ\gamma = 0 \ \text{のとき, }\)\\
\(\todai[n]{\alpha, E^n\beta, E^n\gamma}\ \text{は,}\
  \alpha\circ E^n\pit{i+1}{k} + \pit{j+n+1}{m}\circ E^{n+1}\gamma\
  \text{の coset で, }\pit{i+n+1}{m} \ \text{の部分集合である. }\)

\noindent
\textbf{Prop5.8.}\
\(\pi^{2}_{6}=\mathbb Z_4\gen{\eta_{2}\circ\nu'},\) \quad
\(\pi^{3}_{7}=\mathbb Z_2\gen{\nu'\circ\eta_6},\) \quad
\(\pi^{4}_{8}=\mathbb Z_2\gen{\nu_{4}\circ\eta_7}\oplus \mathbb Z_2\gen{E\nu'\circ\eta_7},\) \\
\(\pi^{5}_{9}=\mathbb Z_2\gen{\nu_{5}\circ\eta_8},\) \quad
\(\pi^{n}_{n+4}=0\ (n\ge6),\) \
\((G_{4};2)=0.\)

\noindent
\textbf{Prop1.3.}\
\(-E\todai[n]{\alpha, E^n\beta, E^n\gamma}
\subset \todai[n+1]{E\alpha, E^{n+1}\beta, E^{n+1}\gamma}\).

\noindent
\textbf{Lem5.13.}\
\(\todai[t]{\nu_5,8\iota_8,\nu_8},\ (0\leq t\leq3)\) 
は, ただ$1$つの元から成る. \\
この元を \(\sigma'''\) と書くと, 
\(H(\sigma''')=4\nu_9=\eta_9^3\) であり, \\
\(\pi_{12}^5=\mathbb{Z}_2\gen{\sigma'''}\)
である. 

\noindent
\textbf{Prop1.2.}\
\textbf{0).}\ $\alpha$, $\beta$, $\gamma$ のうち, いずれかが $0$ なら,
\(
\{\alpha,\ \Ex^n\beta,\ \Ex^n\gamma\}_n\equiv0.
\)\\
\textbf{i).}\ $\alpha\circ \Ex^n\beta=\beta\circ\gamma=0$ のとき,
\(
\{\alpha,\ \Ex^n\beta,\ \Ex^n\gamma\}_n\circ\Ex^{n+1}\delta\ \subset\ \{\alpha,\Ex^n \beta, \Ex^n(\gamma\circ\delta)\}_n.
\)\\
\textbf{ii).}\ $\alpha\circ \Ex^n\beta=\beta\circ\gamma\circ\delta=0$ のとき,
\(
\{\alpha,\ \Ex^n\beta,\ \Ex^n(\gamma\circ\delta)\}_n\ \subset\ \{\alpha,\ \Ex^n(\beta\circ\gamma),\ \Ex^n\delta\}_n .
\)\\
\textbf{iii).}\ $\alpha\circ\Ex^n\beta\circ\Ex^n\gamma=\gamma\circ\delta=0$ なら,
\(
\{\alpha\circ\Ex^n\beta,\ \Ex^n\gamma,\ \Ex^n\delta)\}_n\ \subset\ \{\alpha,\ \Ex^n(\beta\circ\gamma),\ \Ex^n\delta\}_n .
\)\\
\textbf{iv).}\ $\beta\circ\Ex^n\gamma=\gamma\circ\delta=0$ のとき,
\(
\alpha\circ\{\beta,\ \Ex^n\gamma,\ \Ex^n\delta\}_n\ \subset\ \{\alpha\circ\beta,\ \Ex^n\gamma,\ \Ex^n\delta\}_n .
\)

\noindent
\textbf{(5.14).}\
次の列は完全列である.\\
\(
0\to \pi^{5}_{12}\xrightarrow{E}\pi^{6}_{13}\xrightarrow{H}\pi^{11}_{13}\to0,
\quad
0\to \pi^{6}_{13}\xrightarrow{E}\pi^{7}_{14}\xrightarrow{H}\pi^{13}_{14}\to0 .
\)

\noindent
\textbf{(4.8).}\
\( n \ \text{を偶数とする. } \)\\
\(\Delta(\iota_{2n-1})=[\iota_{n-1},\iota_{n-1}]\ne0 \ \text{のとき, }\)\\
\(\pit{2n-1}{n}
=E\,\pit{2n-2}{\,n-1}\ \oplus\ \Delta\,\pit{2n+1}{\,2n+1}
\cong E\,\pit{2n-2}{\,n-1}\ \oplus\ \mathbb{Z},\ \pit{2n}{\,n+1} = E^{2}\,\pit{2n-2}{\,n-1} 
\cong E\,\pit{2n-2}{\,n}.\)\\
\(\Delta(\iota_{2n-1})=0 \ \text{のとき, }\)\\
\(E^{2}:\pit{2n-2}{\,n-1}\to\pit{2n}{\,n+1}\ \text{は単射で, }\
\pit{2n}{\,n+1}/E^{2}\pit{2n-2}{\,n-1}\cong\mathbb{Z}_2 . \)\\
\(\text{さらに,}\ \alpha\in\pit{2n-1}{n} \ \text{に対して, 次が成り立つ. }\)
\(E\alpha\in E^{2}\pit{2n-2}{\,n-1} \ \Leftrightarrow \
H(\alpha)\in 2\,\pit{2n-1}{\,2n-1} .\)



\end{paracol}

\begin{paracol}{2}

\noindent
\(\sigma'\in\pi^{7}_{14}\ ,\ \sigma''\in\pi^{6}_{13}\) 
が一意的に存在して,
\(
E^{2}\sigma'=2x\,E\alpha^{\ast},\ E^{3}\sigma''=4x\,E\alpha^{\ast} .
\)\\
したがって,
\(\pi^{7}_{14}=\gen{\sigma'}\cong\mathbb Z_{8}\ , \
\pi_{13}^6=\gen{\sigma''}\cong\mathbb{Z}_4\) である. \\
また, 明らかに, \(2\sigma'=E\sigma'',\ 2\sigma''=E\sigma'''\) が成り立つ. \\
(5.14) の完全性から, 
\(H(\sigma')=\eta_{13}\), \(H(\sigma'')=\eta_{11}^{2} .\)\\
さらに, \(xE\alpha^{\ast}\notin \im\{E^2:\pi_{10}^3\to\pi_{12}^5\}\) である. \\
(4.8) より, \(x\alpha^{\ast}\) の Hopf invariant は奇数である. \\
Prop2.7 により, 
\(H(\Delta(\iota_{17}))=\pm 2\iota_{15}\) なので, \\
ある整数 \(y\) が存在し, 
\(H(x\alpha^{\ast})=H(y\Delta(\iota_{17}))+\iota_{15} .\)
よって, 
\[
\sigma_{8}:=x\alpha^{\ast}-y\,\Delta(\iota_{17})
\]
と置けば, 
\(H(\sigma_{8})=\iota_{15}\ , \ E\sigma_{8}=xE\alpha^{\ast}\)
である. 
\qed

\(\sigma_n :=E^{n-8}\sigma_8\ , \ (n\geq8) \ , \ \sigma:=E^{\infty}\sigma_8\) と書く. \\
(4.5) より, \(\pi_{n+7}^n \ (n\geq8)\),
\((G_7;2)\) は,
\(\pi_{16}^9\cong\mathbb{Z}_{16}\) に同型で, 
それぞれ,
\(\sigma_n,\ \sigma\) で生成される. \\
これで, \(\pi_{n+7}^n\) は, 
\(\pi_{15}^8\) 以外は全て計算された. \\
\(\pi_{15}^8\) は, Prop4.4 から出る, 
次の (5.15) により計算される. 
\begin{equation}\label{eq:5.15}
% \noindent
% \textbf{(5.15).}\
(\alpha,\beta)\ \longmapsto\ E\alpha+\sigma_{8}\circ\beta
\quad\text{は,}\quad
\pi^{7}_{i-1}\oplus\pi^{15}_{i}\xrightarrow{\cong}\pi^{8}_{i}
\ \text{の同型を与える. }
\end{equation}

\(\pi_{n+7}^n\) に関する結果を, 以下にまとめる. 

\begin{proposition}\label{prop:5.15}
\ \\
% \textbf{Prop5.15.}\
\(\pi_9^2=\pi_{10}^3=\pi_{11}^4=0,\) \\
\(\pi^{5}_{12}=\mathbb Z_2\gen{\sigma'''},\) \\
\(\pi^{6}_{13}=\mathbb Z_4\gen{\sigma''},\) \\
\(\pi^{7}_{14}=\mathbb Z_8\gen{\sigma'},\) \\
\(\pi^{8}_{15}=\mathbb Z\gen{\sigma_8}\oplus\mathbb Z\gen{E\sigma'},\) \\
\(\pi^{n}_{n+7}=\mathbb Z_{16}\gen{\sigma_n}\ (n\ge9),\) \
\((G_{7};2)=\mathbb Z_{16}\gen{\sigma}.\)
\end{proposition}

完全列 
\(\pi^{17}_{17}\xrightarrow{\Delta}\pi^{8}_{15}\xrightarrow{E}\pi^{9}_{16}\)
を考える. \\
\(\Ker E\) は, \(2\sigma_{8}-E\sigma'\) で生成される. 
したがって, 
\begin{equation}\label{eq:5.16}
% \noindent
% \textbf{(5.16).}\
\Delta(\iota_{17})=\pm\,(2\sigma_{8}-E\sigma').
\end{equation}
Lem5.14 の証明の最初と, \(\sigma_8\) の定義から, 次が成り立つ. 
\begin{lemma}\label{lem:5.16}
% \noindent
% \textbf{Lem5.16.}\
\(t>0\) とし, \(\beta\in\pi_{t+4}(S^{m})\) が 
\(\beta\circ \nu_{t+4}=0\) を満たすとする. \\
このとき, ある奇数 \(x\) が存在して,
\[
E^{4}\beta\circ \sigma_{t+8}\ \in\
(-1)^{m}x\todai[7]{\nu_{m+4}, E^{n}\beta, \nu_{t+11}}
+ (-1)^{t}x \todai[t+3]{E^{4}\beta, \nu_{t+8}, 2\nu_{t+11}}.
\]
\end{lemma}

\switchcolumn

\noindent
\textbf{(5.14).}\
次の列は完全列である.\\
\(
0\to \pi^{5}_{12}\xrightarrow{E}\pi^{6}_{13}\xrightarrow{H}\pi^{11}_{13}\to0,
\quad
0\to \pi^{6}_{13}\xrightarrow{E}\pi^{7}_{14}\xrightarrow{H}\pi^{13}_{14}\to0 .
\)

\noindent
\textbf{(4.8).}\
\( n \ \text{を偶数とする. } \)\\
\(\Delta(\iota_{2n-1})=[\iota_{n-1},\iota_{n-1}]\ne0 \ \text{のとき, }\)\\
\(\pit{2n-1}{n}
=E\,\pit{2n-2}{\,n-1}\ \oplus\ \Delta\,\pit{2n+1}{\,2n+1}
\cong E\,\pit{2n-2}{\,n-1}\ \oplus\ \mathbb{Z},\ \pit{2n}{\,n+1} = E^{2}\,\pit{2n-2}{\,n-1} 
\cong E\,\pit{2n-2}{\,n}.\)\\
\(\Delta(\iota_{2n-1})=0 \ \text{のとき, }\)\\
\(E^{2}:\pit{2n-2}{\,n-1}\to\pit{2n}{\,n+1}\ \text{は単射で, }\
\pit{2n}{\,n+1}/E^{2}\pit{2n-2}{\,n-1}\cong\mathbb{Z}_2 . \)\\
\(\text{さらに,}\ \alpha\in\pit{2n-1}{n} \ \text{に対して, 次が成り立つ. }\)\\
\(E\alpha\in E^{2}\pit{2n-2}{\,n-1} \ \Leftrightarrow \
H(\alpha)\in 2\,\pit{2n-1}{\,2n-1} .\)

\noindent
\textbf{Prop2.7.}\
\(H\Delta(\iota_{2n+1}) = \pm2\iota_{2m-1},\ (n:\text{even})\).

\noindent
\textbf{(4.5).}\
\(E^{m-n}:\pit{n+k}{n}\xrightarrow{\cong}\pit{m+k}{m},\
(n\ge k+2). \)

\noindent
\textbf{Prop4.4.}\
\(H(\alpha) = \pm 2\iota_{2m-1}\ \Longrightarrow\
\pit{i-1}{m-1}\oplus\pit{i}{2m-1}\xrightarrow{\cong}\pit{i}{m},\
(\beta,\gamma)\mapsto E\beta+\alpha\circ\gamma.\)\\
\(0=\pit{2}{1}\to\pit{3}{2}
\xrightarrow{H}{\cong}\pit{3}{3}
\xrightarrow{\Delta}\pit{1}{1}
\xrightarrow{E}\pit{2}{2}.\)


\noindent
Th 3.6 を \(\alpha=\nu_{4}\) に適用し, \\
Prop \ref{prop:5.11} により, 
\(\nu_4\) は (3.1) の条件を満たす. \\
\(
2E^{4-3}\nu_{4}\circ E^{7-3}\nu_{4}=2\nu_{5}^{2}=0,
\quad (h=3).
\)\\
このとき, 
\(\alpha^{\ast}\in\pi^{8}_{15}\) が存在し, 
\(\beta\circ\nu_{t+4}=0\) を満たす任意の \(\beta\in\pi_{t+4}(S^m)\) に対して, \\
\(
E^{4}\beta\circ E^t\alpha^{\ast}
\in (-1)^{m}\todai[7]{\nu_{m+4},\,E^{7}\beta,\,\nu_{t+11}}
+(-1)^{t}\todai[t+3]{E^{4}\beta,\,\nu_{t+8},\,2\nu_{t+11}} .
\)





\end{paracol}

\newpage

\begin{paracol}{2}

\renewcommand{\refname}{References}
\begin{thebibliography}{99}

\bibitem{Adams1960}
J.~F.~Adams,
\newblock On the non-existence of elements of Hopf invariant one,
\newblock Ann. of Math., {\bf 72} (1960), 20--104.

\bibitem{Adem}
J.~Adem,
\newblock The relations on Steenrod powers of cohomology classes,
\newblock Algebraic geometry and topology, Princeton.

\bibitem{Barratt1955}
M.~G.~Barratt,
\newblock Track groups I, II,
\newblock Proc. London Math. Soc., {\bf 5} (1955), 71--106, 285--329.

\bibitem{Barratt1961}
M.~G.~Barratt,
\newblock Note on a formula due to Toda,
\newblock Jour. London Math. Soc., {\bf 36} (1961), 95--96.

\bibitem{BarrattHilton1953}
M.~G.~Barratt and P.~J.~Hilton,
\newblock On join operations of homotopy groups,
\newblock Proc. London Math. Soc., {\bf 3} (1953), 430--445.

\bibitem{BlakersMassey1952}
A.~L.~Blakers and W.~S.~Massey,
\newblock The homotopy groups of a triad II,
\newblock Ann. of Math., {\bf 55} (1952), 192--201.

\bibitem{Freudenthal1937}
H.~Freudenthal,
\newblock Über die Klassen der Sphärenabbildungen,
\newblock Compos. Math., {\bf 5} (1937).

\bibitem{Hopf1935}
H.~Hopf,
\newblock Über die Abbildungen von Sphären auf Sphären niedriger Dimension,
\newblock Fund. Math., {\bf 25} (1935).

\bibitem{James1955}
I.~M.~James,
\newblock Reduced product spaces,
\newblock Ann. of Math., {\bf 62} (1955), 170--197.

\bibitem{James1956}
I.~M.~James,
\newblock Suspension triad of a sphere,
\newblock Ann. of Math., {\bf 63} (1956), 407--429.

\bibitem{James1958}
I.~M.~James,
\newblock The intrinsic join: A study of the homotopy groups of Stiefel manifold,
\newblock Proc. London Math. Soc., {\bf 8} (1958), 507--535.

\bibitem{Moore1953}
J.~C.~Moore,
\newblock Some applications of homology theory to homotopy problems,
\newblock Ann. of Math., {\bf 58} (1953), 325--350.

\bibitem{Serre1951}
J.-P.~Serre,
\newblock Homologie singulière des espaces fibrés,
\newblock Ann. of Math., {\bf 54} (1951), 425--505.

\bibitem{Serre1953a}
J.-P.~Serre,
\newblock Cohomologie modulo 2 des complexes d'Eilenberg-MacLane,
\newblock Comm. Math. Helv., {\bf 27} (1953), 198--231.

\bibitem{Serre1953b}
J.-P.~Serre,
\newblock Groups d'homotopie et classes groupes abéliens,
\newblock Ann. of Math., {\bf 58} (1953), 258--294.

\bibitem{Serre1953c}
J.-P.~Serre,
\newblock Quelques calculs de groupes d'homotopie,
\newblock Comptes Rendus Paris, {\bf 236} (1953), 2475--2477.

\bibitem{Steenrod1951}
N.~E.~Steenrod,
\newblock Topology of fibre bundles,
\newblock Princeton.

\bibitem{Steenrod1949}
N.~E.~Steenrod,
\newblock Cohomology invariants of mappings,
\newblock Ann. of Math., {\bf 50} (1949), 954--988.

\bibitem{Steenrod1953}
N.~E.~Steenrod,
\newblock Cyclic reduced powers of cohomology classes,
\newblock Proc. Nat. Acad. Sci. U.S.A., {\bf 39} (1953), 217--223.

\bibitem{Toda1952}
H.~Toda,
\newblock Generalized Whitehead products and homotopy groups of spheres,
\newblock Jour. of Inst. Poly. Osaka City Univ., {\bf 3} (1952), 43--82.

\bibitem{Toda1956}
H.~Toda,
\newblock On double suspension $E^2$,
\newblock Jour. of Inst. Poly. Osaka City Univ., {\bf 7} (1956), 103--145.

\bibitem{Toda1957}
H.~Toda,
\newblock Reduced join and Whitehead product,
\newblock Jour. of Inst. Poly. Osaka City Univ., {\bf 8} (1957), 15--30.

\bibitem{Toda1958a}
H.~Toda,
\newblock On exact sequences in Steenrod algebra mod 2,
\newblock Memoirs Univ. of Kyoto, {\bf 31} (1958), 33--64.

\bibitem{Toda1958b}
H.~Toda,
\newblock $p$-primary components of homotopy groups II. mod $p$ Hopf invariant,
\newblock Memoirs Univ. of Kyoto, {\bf 31} (1958), 143--160.

\bibitem{Toda1959}
H.~Toda,
\newblock $p$-primary components of homotopy groups IV. Compositions and toric constructions,
\newblock Memoirs Univ. of Kyoto, {\bf 32} (1959), 297--332.

\bibitem{Whitehead1950a}
G.~W.~Whitehead,
\newblock A generalization of the Hopf invariant,
\newblock Ann. of Math., {\bf 51} (1950), 192--237.

\bibitem{Whitehead1950b}
G.~W.~Whitehead,
\newblock The $(n+2)$nd homotopy group of the $n$-sphere,
\newblock Ann. of Math., {\bf 52} (1950), 245--247.

\bibitem{Whitehead1941}
J.~H.~C.~Whitehead,
\newblock On adding relations to homotopy groups,
\newblock Ann. of Math., {\bf 42} (1941), 409--428.

\bibitem{Whitehead1944}
J.~H.~C.~Whitehead,
\newblock On the groups $\pi_n(V_{m,n})$ and sphere bundles,
\newblock Proc. London Math. Soc., {\bf 48} (1944), 243--291.

\end{thebibliography}

\end{paracol}


\end{document}
